% paper65-ci-cd-integration.tex
% Integrating JuGeo into CI/CD Pipelines for Continuous Formal Verification
% JuGeo Paper Series — Paper 65

\documentclass[11pt]{article}
% jugeo-common.tex — shared preamble for all Judgment Geometry papers
% Include via: % jugeo-common.tex — shared preamble for all Judgment Geometry papers
% Include via: % jugeo-common.tex — shared preamble for all Judgment Geometry papers
% Include via: \input{jugeo-common}

% ── Packages ────────────────────────────────────────────────────────────
\usepackage{amsmath,amssymb,amsthm,mathtools}
\usepackage{tikz-cd}
\usepackage{listings}
\usepackage{xcolor}
\usepackage{xspace}
\usepackage{booktabs}
\usepackage{multirow}
\usepackage{enumitem}
\usepackage[expansion=false]{microtype}
\usepackage{hyperref}
\usepackage{cleveref}
% \usepackage{thmtools}  % disabled: not available in this TeX installation
\usepackage{float}
\usepackage{caption}
\usepackage{subcaption}
\usepackage{tabularx}
\usepackage{stmaryrd}       % \llbracket, \rrbracket
\usepackage{bbm}            % \mathbbm{1}

% ── Page geometry ───────────────────────────────────────────────────────
\usepackage[margin=1in]{geometry}

% ── Theorem environments ────────────────────────────────────────────────
\theoremstyle{plain}
\newtheorem{theorem}{Theorem}[section]
\newtheorem{lemma}[theorem]{Lemma}
\newtheorem{proposition}[theorem]{Proposition}
\newtheorem{corollary}[theorem]{Corollary}

\theoremstyle{definition}
\newtheorem{definition}[theorem]{Definition}
\newtheorem{example}[theorem]{Example}
\newtheorem{construction}[theorem]{Construction}

\theoremstyle{remark}
\newtheorem{remark}[theorem]{Remark}
\newtheorem{notation}[theorem]{Notation}

% ── Python listings style ──────────────────────────────────────────────
\definecolor{jugeoBlue}{HTML}{2563EB}
\definecolor{jugeoGreen}{HTML}{059669}
\definecolor{jugeoRed}{HTML}{DC2626}
\definecolor{jugeoPurple}{HTML}{7C3AED}
\definecolor{jugeoGray}{HTML}{6B7280}
\definecolor{jugeoBg}{HTML}{F8FAFC}

\lstdefinestyle{jugeo-python}{
  language=Python,
  basicstyle=\ttfamily\small,
  keywordstyle=\color{jugeoBlue}\bfseries,
  stringstyle=\color{jugeoGreen},
  commentstyle=\color{jugeoGray}\itshape,
  numberstyle=\tiny\color{jugeoGray},
  backgroundcolor=\color{jugeoBg},
  numbers=left,
  numbersep=6pt,
  frame=single,
  framerule=0.4pt,
  rulecolor=\color{jugeoGray!40},
  breaklines=true,
  breakatwhitespace=true,
  showstringspaces=false,
  tabsize=4,
  xleftmargin=1.5em,
  framexleftmargin=1.5em,
  morekeywords={dataclass,frozen,field,Enum,IntEnum,auto,Any,
                Mapping,Sequence,Optional,match,case},
  literate={→}{{$\to$}}1 {←}{{$\leftarrow$}}1
           {≤}{{$\leq$}}1 {≥}{{$\geq$}}1 {≠}{{$\neq$}}1
           {∈}{{$\in$}}1 {∀}{{$\forall$}}1 {∃}{{$\exists$}}1
           {⊕}{{$\oplus$}}1 {⊖}{{$\ominus$}}1,
  escapeinside={(*@}{@*)},
}
\lstset{style=jugeo-python}

\lstdefinestyle{jugeo-output}{
  basicstyle=\ttfamily\small,
  backgroundcolor=\color{jugeoBg},
  frame=single,
  framerule=0.4pt,
  rulecolor=\color{jugeoGray!40},
  breaklines=true,
  showstringspaces=false,
  xleftmargin=1.5em,
  framexleftmargin=0.5em,
  numbers=none,
}

% ── JuGeo-specific macros ──────────────────────────────────────────────

% Judgment 8-tuple
\newcommand{\J}{\mathcal{J}}
\newcommand{\Jud}[8]{(#1,\,#2,\,#3,\,#4,\,#5,\,#6,\,#7,\,#8)}
\newcommand{\JudShort}[1]{\J_{#1}}

% Components
\newcommand{\coord}{\mathsf{c}}            % coordinate
\newcommand{\prop}{\varphi}                 % proposition
\newcommand{\carrier}{\mathcal{A}}          % carrier type
\newcommand{\evid}{\mathcal{E}}             % evidence bundle
\newcommand{\oblig}{\mathcal{O}}            % obligations
\newcommand{\obstr}{\mathcal{B}}            % obstructions
\newcommand{\trust}{\mathcal{T}}            % trust annotation
\newcommand{\prov}{\Pi}                     % provenance

% Trust algebra
\newcommand{\Talg}{\mathfrak{T}}            % trust ordered algebra
\newcommand{\Eadm}{\mathcal{E}_{\mathrm{adm}}}
\newcommand{\tleq}{\preceq}                 % trust ordering
\newcommand{\tjoin}{\oplus}                 % conservative join
\newcommand{\tatten}{\ominus}               % attenuation
\newcommand{\tpromote}{\uparrow_\pi}        % promotion
\newcommand{\tdemote}{\downarrow_\chi}       % demotion

% Trust tiers
\newcommand{\Tcontra}{\ensuremath{\mathsf{CONTRADICTED}}}
\newcommand{\Tunver}{\ensuremath{\mathsf{UNVERIFIED}}}
\newcommand{\Tcopilot}{\ensuremath{\mathsf{COPILOT}}}
\newcommand{\Toracle}{\ensuremath{\mathsf{ORACLE}}}
\newcommand{\Truntime}{\ensuremath{\mathsf{RUNTIME}}}
\newcommand{\Tsolver}{\ensuremath{\mathsf{SOLVER}}}
\newcommand{\Tproof}{\ensuremath{\mathsf{PROOF}}}

% Site and geometry
\newcommand{\Site}{\mathcal{S}}             % semantic site
\newcommand{\Cov}{\mathrm{Cov}}            % covering family
\newcommand{\Ob}{\mathrm{Ob}}              % objects
\newcommand{\Mor}{\mathrm{Mor}}            % morphisms
\newcommand{\res}{\mathrm{res}}            % restriction
\newcommand{\Cech}{\check{C}}              % Čech complex
\newcommand{\Hcoh}[1]{H^{#1}}             % cohomology group

% Descent
\newcommand{\descent}{\mathrm{desc}}
\newcommand{\glue}{\mathrm{glue}}
\newcommand{\overlap}[2]{#1 \cap #2}

% Semantic moves
\newcommand{\VERIFY}{\textsc{Verify}}
\newcommand{\CONSTRUCT}{\textsc{Construct}}
\newcommand{\NORMALIZE}{\textsc{Normalize}}
\newcommand{\GLUE}{\textsc{Glue}}
\newcommand{\RESTRICT}{\textsc{Restrict}}
\newcommand{\TRANSPORT}{\textsc{Transport}}
\newcommand{\REPAIR}{\textsc{Repair}}
\newcommand{\PROMOTE}{\textsc{Promote}}
\newcommand{\CHALLENGE}{\textsc{Challenge}}
\newcommand{\ESCALATE}{\textsc{Escalate}}

% Fragments
\newcommand{\QFLIA}{\texttt{QF\_LIA}}
\newcommand{\QFLRA}{\texttt{QF\_LRA}}
\newcommand{\QFBV}{\texttt{QF\_BV}}
\newcommand{\QFUF}{\texttt{QF\_UF}}

% Misc
\newcommand{\Python}{\textsc{Python}}
\newcommand{\JuGeo}{\textsc{JuGeo}}
\newcommand{\LEAN}{\textsc{Lean}\,4}
\newcommand{\Fstar}{F$^{\star}$}
\newcommand{\Coq}{\textsc{Coq}}
\newcommand{\Dafny}{\textsc{Dafny}}

% Common shortcuts
\newcommand{\ie}{i.e.\@\xspace}
\newcommand{\eg}{e.g.\@\xspace}
\newcommand{\cf}{cf.\@\xspace}
\newcommand{\wrt}{w.r.t.\@\xspace}

% Evaluation shortcuts
\newcommand{\numcases}{300}
\newcommand{\numspec}{100}
\newcommand{\numeq}{100}
\newcommand{\numbug}{100}
\newcommand{\medtime}{6.5\,ms}
\newcommand{\totaltime}{1.949\,s}


% ── Packages ────────────────────────────────────────────────────────────
\usepackage{amsmath,amssymb,amsthm,mathtools}
\usepackage{tikz-cd}
\usepackage{listings}
\usepackage{xcolor}
\usepackage{xspace}
\usepackage{booktabs}
\usepackage{multirow}
\usepackage{enumitem}
\usepackage[expansion=false]{microtype}
\usepackage{hyperref}
\usepackage{cleveref}
% \usepackage{thmtools}  % disabled: not available in this TeX installation
\usepackage{float}
\usepackage{caption}
\usepackage{subcaption}
\usepackage{tabularx}
\usepackage{stmaryrd}       % \llbracket, \rrbracket
\usepackage{bbm}            % \mathbbm{1}

% ── Page geometry ───────────────────────────────────────────────────────
\usepackage[margin=1in]{geometry}

% ── Theorem environments ────────────────────────────────────────────────
\theoremstyle{plain}
\newtheorem{theorem}{Theorem}[section]
\newtheorem{lemma}[theorem]{Lemma}
\newtheorem{proposition}[theorem]{Proposition}
\newtheorem{corollary}[theorem]{Corollary}

\theoremstyle{definition}
\newtheorem{definition}[theorem]{Definition}
\newtheorem{example}[theorem]{Example}
\newtheorem{construction}[theorem]{Construction}

\theoremstyle{remark}
\newtheorem{remark}[theorem]{Remark}
\newtheorem{notation}[theorem]{Notation}

% ── Python listings style ──────────────────────────────────────────────
\definecolor{jugeoBlue}{HTML}{2563EB}
\definecolor{jugeoGreen}{HTML}{059669}
\definecolor{jugeoRed}{HTML}{DC2626}
\definecolor{jugeoPurple}{HTML}{7C3AED}
\definecolor{jugeoGray}{HTML}{6B7280}
\definecolor{jugeoBg}{HTML}{F8FAFC}

\lstdefinestyle{jugeo-python}{
  language=Python,
  basicstyle=\ttfamily\small,
  keywordstyle=\color{jugeoBlue}\bfseries,
  stringstyle=\color{jugeoGreen},
  commentstyle=\color{jugeoGray}\itshape,
  numberstyle=\tiny\color{jugeoGray},
  backgroundcolor=\color{jugeoBg},
  numbers=left,
  numbersep=6pt,
  frame=single,
  framerule=0.4pt,
  rulecolor=\color{jugeoGray!40},
  breaklines=true,
  breakatwhitespace=true,
  showstringspaces=false,
  tabsize=4,
  xleftmargin=1.5em,
  framexleftmargin=1.5em,
  morekeywords={dataclass,frozen,field,Enum,IntEnum,auto,Any,
                Mapping,Sequence,Optional,match,case},
  literate={→}{{$\to$}}1 {←}{{$\leftarrow$}}1
           {≤}{{$\leq$}}1 {≥}{{$\geq$}}1 {≠}{{$\neq$}}1
           {∈}{{$\in$}}1 {∀}{{$\forall$}}1 {∃}{{$\exists$}}1
           {⊕}{{$\oplus$}}1 {⊖}{{$\ominus$}}1,
  escapeinside={(*@}{@*)},
}
\lstset{style=jugeo-python}

\lstdefinestyle{jugeo-output}{
  basicstyle=\ttfamily\small,
  backgroundcolor=\color{jugeoBg},
  frame=single,
  framerule=0.4pt,
  rulecolor=\color{jugeoGray!40},
  breaklines=true,
  showstringspaces=false,
  xleftmargin=1.5em,
  framexleftmargin=0.5em,
  numbers=none,
}

% ── JuGeo-specific macros ──────────────────────────────────────────────

% Judgment 8-tuple
\newcommand{\J}{\mathcal{J}}
\newcommand{\Jud}[8]{(#1,\,#2,\,#3,\,#4,\,#5,\,#6,\,#7,\,#8)}
\newcommand{\JudShort}[1]{\J_{#1}}

% Components
\newcommand{\coord}{\mathsf{c}}            % coordinate
\newcommand{\prop}{\varphi}                 % proposition
\newcommand{\carrier}{\mathcal{A}}          % carrier type
\newcommand{\evid}{\mathcal{E}}             % evidence bundle
\newcommand{\oblig}{\mathcal{O}}            % obligations
\newcommand{\obstr}{\mathcal{B}}            % obstructions
\newcommand{\trust}{\mathcal{T}}            % trust annotation
\newcommand{\prov}{\Pi}                     % provenance

% Trust algebra
\newcommand{\Talg}{\mathfrak{T}}            % trust ordered algebra
\newcommand{\Eadm}{\mathcal{E}_{\mathrm{adm}}}
\newcommand{\tleq}{\preceq}                 % trust ordering
\newcommand{\tjoin}{\oplus}                 % conservative join
\newcommand{\tatten}{\ominus}               % attenuation
\newcommand{\tpromote}{\uparrow_\pi}        % promotion
\newcommand{\tdemote}{\downarrow_\chi}       % demotion

% Trust tiers
\newcommand{\Tcontra}{\ensuremath{\mathsf{CONTRADICTED}}}
\newcommand{\Tunver}{\ensuremath{\mathsf{UNVERIFIED}}}
\newcommand{\Tcopilot}{\ensuremath{\mathsf{COPILOT}}}
\newcommand{\Toracle}{\ensuremath{\mathsf{ORACLE}}}
\newcommand{\Truntime}{\ensuremath{\mathsf{RUNTIME}}}
\newcommand{\Tsolver}{\ensuremath{\mathsf{SOLVER}}}
\newcommand{\Tproof}{\ensuremath{\mathsf{PROOF}}}

% Site and geometry
\newcommand{\Site}{\mathcal{S}}             % semantic site
\newcommand{\Cov}{\mathrm{Cov}}            % covering family
\newcommand{\Ob}{\mathrm{Ob}}              % objects
\newcommand{\Mor}{\mathrm{Mor}}            % morphisms
\newcommand{\res}{\mathrm{res}}            % restriction
\newcommand{\Cech}{\check{C}}              % Čech complex
\newcommand{\Hcoh}[1]{H^{#1}}             % cohomology group

% Descent
\newcommand{\descent}{\mathrm{desc}}
\newcommand{\glue}{\mathrm{glue}}
\newcommand{\overlap}[2]{#1 \cap #2}

% Semantic moves
\newcommand{\VERIFY}{\textsc{Verify}}
\newcommand{\CONSTRUCT}{\textsc{Construct}}
\newcommand{\NORMALIZE}{\textsc{Normalize}}
\newcommand{\GLUE}{\textsc{Glue}}
\newcommand{\RESTRICT}{\textsc{Restrict}}
\newcommand{\TRANSPORT}{\textsc{Transport}}
\newcommand{\REPAIR}{\textsc{Repair}}
\newcommand{\PROMOTE}{\textsc{Promote}}
\newcommand{\CHALLENGE}{\textsc{Challenge}}
\newcommand{\ESCALATE}{\textsc{Escalate}}

% Fragments
\newcommand{\QFLIA}{\texttt{QF\_LIA}}
\newcommand{\QFLRA}{\texttt{QF\_LRA}}
\newcommand{\QFBV}{\texttt{QF\_BV}}
\newcommand{\QFUF}{\texttt{QF\_UF}}

% Misc
\newcommand{\Python}{\textsc{Python}}
\newcommand{\JuGeo}{\textsc{JuGeo}}
\newcommand{\LEAN}{\textsc{Lean}\,4}
\newcommand{\Fstar}{F$^{\star}$}
\newcommand{\Coq}{\textsc{Coq}}
\newcommand{\Dafny}{\textsc{Dafny}}

% Common shortcuts
\newcommand{\ie}{i.e.\@\xspace}
\newcommand{\eg}{e.g.\@\xspace}
\newcommand{\cf}{cf.\@\xspace}
\newcommand{\wrt}{w.r.t.\@\xspace}

% Evaluation shortcuts
\newcommand{\numcases}{300}
\newcommand{\numspec}{100}
\newcommand{\numeq}{100}
\newcommand{\numbug}{100}
\newcommand{\medtime}{6.5\,ms}
\newcommand{\totaltime}{1.949\,s}


% ── Packages ────────────────────────────────────────────────────────────
\usepackage{amsmath,amssymb,amsthm,mathtools}
\usepackage{tikz-cd}
\usepackage{listings}
\usepackage{xcolor}
\usepackage{xspace}
\usepackage{booktabs}
\usepackage{multirow}
\usepackage{enumitem}
\usepackage[expansion=false]{microtype}
\usepackage{hyperref}
\usepackage{cleveref}
% \usepackage{thmtools}  % disabled: not available in this TeX installation
\usepackage{float}
\usepackage{caption}
\usepackage{subcaption}
\usepackage{tabularx}
\usepackage{stmaryrd}       % \llbracket, \rrbracket
\usepackage{bbm}            % \mathbbm{1}

% ── Page geometry ───────────────────────────────────────────────────────
\usepackage[margin=1in]{geometry}

% ── Theorem environments ────────────────────────────────────────────────
\theoremstyle{plain}
\newtheorem{theorem}{Theorem}[section]
\newtheorem{lemma}[theorem]{Lemma}
\newtheorem{proposition}[theorem]{Proposition}
\newtheorem{corollary}[theorem]{Corollary}

\theoremstyle{definition}
\newtheorem{definition}[theorem]{Definition}
\newtheorem{example}[theorem]{Example}
\newtheorem{construction}[theorem]{Construction}

\theoremstyle{remark}
\newtheorem{remark}[theorem]{Remark}
\newtheorem{notation}[theorem]{Notation}

% ── Python listings style ──────────────────────────────────────────────
\definecolor{jugeoBlue}{HTML}{2563EB}
\definecolor{jugeoGreen}{HTML}{059669}
\definecolor{jugeoRed}{HTML}{DC2626}
\definecolor{jugeoPurple}{HTML}{7C3AED}
\definecolor{jugeoGray}{HTML}{6B7280}
\definecolor{jugeoBg}{HTML}{F8FAFC}

\lstdefinestyle{jugeo-python}{
  language=Python,
  basicstyle=\ttfamily\small,
  keywordstyle=\color{jugeoBlue}\bfseries,
  stringstyle=\color{jugeoGreen},
  commentstyle=\color{jugeoGray}\itshape,
  numberstyle=\tiny\color{jugeoGray},
  backgroundcolor=\color{jugeoBg},
  numbers=left,
  numbersep=6pt,
  frame=single,
  framerule=0.4pt,
  rulecolor=\color{jugeoGray!40},
  breaklines=true,
  breakatwhitespace=true,
  showstringspaces=false,
  tabsize=4,
  xleftmargin=1.5em,
  framexleftmargin=1.5em,
  morekeywords={dataclass,frozen,field,Enum,IntEnum,auto,Any,
                Mapping,Sequence,Optional,match,case},
  literate={→}{{$\to$}}1 {←}{{$\leftarrow$}}1
           {≤}{{$\leq$}}1 {≥}{{$\geq$}}1 {≠}{{$\neq$}}1
           {∈}{{$\in$}}1 {∀}{{$\forall$}}1 {∃}{{$\exists$}}1
           {⊕}{{$\oplus$}}1 {⊖}{{$\ominus$}}1,
  escapeinside={(*@}{@*)},
}
\lstset{style=jugeo-python}

\lstdefinestyle{jugeo-output}{
  basicstyle=\ttfamily\small,
  backgroundcolor=\color{jugeoBg},
  frame=single,
  framerule=0.4pt,
  rulecolor=\color{jugeoGray!40},
  breaklines=true,
  showstringspaces=false,
  xleftmargin=1.5em,
  framexleftmargin=0.5em,
  numbers=none,
}

% ── JuGeo-specific macros ──────────────────────────────────────────────

% Judgment 8-tuple
\newcommand{\J}{\mathcal{J}}
\newcommand{\Jud}[8]{(#1,\,#2,\,#3,\,#4,\,#5,\,#6,\,#7,\,#8)}
\newcommand{\JudShort}[1]{\J_{#1}}

% Components
\newcommand{\coord}{\mathsf{c}}            % coordinate
\newcommand{\prop}{\varphi}                 % proposition
\newcommand{\carrier}{\mathcal{A}}          % carrier type
\newcommand{\evid}{\mathcal{E}}             % evidence bundle
\newcommand{\oblig}{\mathcal{O}}            % obligations
\newcommand{\obstr}{\mathcal{B}}            % obstructions
\newcommand{\trust}{\mathcal{T}}            % trust annotation
\newcommand{\prov}{\Pi}                     % provenance

% Trust algebra
\newcommand{\Talg}{\mathfrak{T}}            % trust ordered algebra
\newcommand{\Eadm}{\mathcal{E}_{\mathrm{adm}}}
\newcommand{\tleq}{\preceq}                 % trust ordering
\newcommand{\tjoin}{\oplus}                 % conservative join
\newcommand{\tatten}{\ominus}               % attenuation
\newcommand{\tpromote}{\uparrow_\pi}        % promotion
\newcommand{\tdemote}{\downarrow_\chi}       % demotion

% Trust tiers
\newcommand{\Tcontra}{\ensuremath{\mathsf{CONTRADICTED}}}
\newcommand{\Tunver}{\ensuremath{\mathsf{UNVERIFIED}}}
\newcommand{\Tcopilot}{\ensuremath{\mathsf{COPILOT}}}
\newcommand{\Toracle}{\ensuremath{\mathsf{ORACLE}}}
\newcommand{\Truntime}{\ensuremath{\mathsf{RUNTIME}}}
\newcommand{\Tsolver}{\ensuremath{\mathsf{SOLVER}}}
\newcommand{\Tproof}{\ensuremath{\mathsf{PROOF}}}

% Site and geometry
\newcommand{\Site}{\mathcal{S}}             % semantic site
\newcommand{\Cov}{\mathrm{Cov}}            % covering family
\newcommand{\Ob}{\mathrm{Ob}}              % objects
\newcommand{\Mor}{\mathrm{Mor}}            % morphisms
\newcommand{\res}{\mathrm{res}}            % restriction
\newcommand{\Cech}{\check{C}}              % Čech complex
\newcommand{\Hcoh}[1]{H^{#1}}             % cohomology group

% Descent
\newcommand{\descent}{\mathrm{desc}}
\newcommand{\glue}{\mathrm{glue}}
\newcommand{\overlap}[2]{#1 \cap #2}

% Semantic moves
\newcommand{\VERIFY}{\textsc{Verify}}
\newcommand{\CONSTRUCT}{\textsc{Construct}}
\newcommand{\NORMALIZE}{\textsc{Normalize}}
\newcommand{\GLUE}{\textsc{Glue}}
\newcommand{\RESTRICT}{\textsc{Restrict}}
\newcommand{\TRANSPORT}{\textsc{Transport}}
\newcommand{\REPAIR}{\textsc{Repair}}
\newcommand{\PROMOTE}{\textsc{Promote}}
\newcommand{\CHALLENGE}{\textsc{Challenge}}
\newcommand{\ESCALATE}{\textsc{Escalate}}

% Fragments
\newcommand{\QFLIA}{\texttt{QF\_LIA}}
\newcommand{\QFLRA}{\texttt{QF\_LRA}}
\newcommand{\QFBV}{\texttt{QF\_BV}}
\newcommand{\QFUF}{\texttt{QF\_UF}}

% Misc
\newcommand{\Python}{\textsc{Python}}
\newcommand{\JuGeo}{\textsc{JuGeo}}
\newcommand{\LEAN}{\textsc{Lean}\,4}
\newcommand{\Fstar}{F$^{\star}$}
\newcommand{\Coq}{\textsc{Coq}}
\newcommand{\Dafny}{\textsc{Dafny}}

% Common shortcuts
\newcommand{\ie}{i.e.\@\xspace}
\newcommand{\eg}{e.g.\@\xspace}
\newcommand{\cf}{cf.\@\xspace}
\newcommand{\wrt}{w.r.t.\@\xspace}

% Evaluation shortcuts
\newcommand{\numcases}{300}
\newcommand{\numspec}{100}
\newcommand{\numeq}{100}
\newcommand{\numbug}{100}
\newcommand{\medtime}{6.5\,ms}
\newcommand{\totaltime}{1.949\,s}

% experiment-data.tex — AUTO-GENERATED from real jugeo CLI runs
% DO NOT EDIT — regenerate with: python3 experiments/run_universal_metrics.py
% Generated from 30 programs across 6 domains

\newcommand{\expTotalPrograms}{30}
\newcommand{\expVerified}{30}
\newcommand{\expAccuracy}{100.0\%}
\newcommand{\expCoordsMin}{2}
\newcommand{\expCoordsMax}{10}
\newcommand{\expCoordsMean}{5.2}
\newcommand{\expCoordsMedian}{5.5}
\newcommand{\expPropsMin}{4}
\newcommand{\expPropsMax}{20}
\newcommand{\expPropsMean}{10.4}
\newcommand{\expPropsSum}{311}
\newcommand{\expPropsOkSum}{310}
\newcommand{\expTimeMin}{3.506\,s}
\newcommand{\expTimeMax}{6.714\,s}
\newcommand{\expTimeMean}{4.64\,s}
\newcommand{\expTimeMedian}{4.508\,s}
\newcommand{\expTimeTotal}{139.204\,s}
\newcommand{\expObstructionTotal}{0}
\newcommand{\expHOneZero}{30}
\newcommand{\expDomainCount}{6}
\newcommand{\expTrustCOPILOTSUGGESTED}{30}
\newcommand{\expDomainDsCount}{5}
\newcommand{\expDomainDsVerified}{5}
\newcommand{\expDomainDsAvgCoords}{7.6}
\newcommand{\expDomainDsAvgProps}{15.2}
\newcommand{\expDomainMathCount}{5}
\newcommand{\expDomainMathVerified}{5}
\newcommand{\expDomainMathAvgCoords}{4.8}
\newcommand{\expDomainMathAvgProps}{9.4}
\newcommand{\expDomainSmCount}{5}
\newcommand{\expDomainSmVerified}{5}
\newcommand{\expDomainSmAvgCoords}{7.6}
\newcommand{\expDomainSmAvgProps}{15.2}
\newcommand{\expDomainSortCount}{5}
\newcommand{\expDomainSortVerified}{5}
\newcommand{\expDomainSortAvgCoords}{2.4}
\newcommand{\expDomainSortAvgProps}{4.8}
\newcommand{\expDomainStrCount}{5}
\newcommand{\expDomainStrVerified}{5}
\newcommand{\expDomainStrAvgCoords}{3.2}
\newcommand{\expDomainStrAvgProps}{6.4}
\newcommand{\expDomainWebCount}{5}
\newcommand{\expDomainWebVerified}{5}
\newcommand{\expDomainWebAvgCoords}{5.6}
\newcommand{\expDomainWebAvgProps}{11.2}

% subsystem-data.tex — AUTO-GENERATED from real jugeo subsystem experiments
% DO NOT EDIT — regenerate with: python3 experiments/run_subsystem_experiments.py

\newcommand{\subCliTotal}{10}
\newcommand{\subCliVerified}{9}
\newcommand{\subCliAccuracy}{90.0\%}
\newcommand{\subCliMeanTime}{4.132\,s}
\newcommand{\subCliMinTime}{3.472\,s}
\newcommand{\subCliMaxTime}{5.372\,s}
\newcommand{\subCliTotalCoords}{54}
\newcommand{\subCliTotalProps}{107}
\newcommand{\subCliTotalPropsOk}{106}
\newcommand{\subSiteCalls}{90}
\newcommand{\subSiteSuccessful}{69}
\newcommand{\subSiteSuccessRate}{76.7\%}
\newcommand{\subSiteMeanTime}{0.0001\,s}
\newcommand{\subSiteTotalTime}{0.005\,s}
\newcommand{\subSpecificationsatisfactionSmallTime}{0.0003\,s}
\newcommand{\subSpecificationsatisfactionMediumTime}{0.0001\,s}
\newcommand{\subSpecificationsatisfactionLargeTime}{0.0001\,s}
\newcommand{\subInhabitantfleetSmallTime}{0.0\,s}
\newcommand{\subInhabitantfleetMediumTime}{0.0\,s}
\newcommand{\subInhabitantfleetLargeTime}{0.0\,s}
\newcommand{\subTheoremecologySmallTime}{0.0\,s}
\newcommand{\subTheoremecologyMediumTime}{0.0\,s}
\newcommand{\subTheoremecologyLargeTime}{0.0\,s}
\newcommand{\subStatespaceexplorationSmallTime}{0.0\,s}
\newcommand{\subStatespaceexplorationMediumTime}{0.0\,s}
\newcommand{\subStatespaceexplorationLargeTime}{0.0\,s}
\newcommand{\subRepairsemanticsSmallTime}{0.0\,s}
\newcommand{\subRepairsemanticsMediumTime}{0.0\,s}
\newcommand{\subRepairsemanticsLargeTime}{0.0\,s}
\newcommand{\subReplaygluingSmallTime}{0.0\,s}
\newcommand{\subReplaygluingMediumTime}{0.0\,s}
\newcommand{\subReplaygluingLargeTime}{0.0\,s}
\newcommand{\subSemanticfuturesSmallTime}{0.0\,s}
\newcommand{\subSemanticfuturesMediumTime}{0.0\,s}
\newcommand{\subSemanticfuturesLargeTime}{0.0\,s}
\newcommand{\subHypercovertreatySmallTime}{0.0\,s}
\newcommand{\subHypercovertreatyMediumTime}{0.0\,s}
\newcommand{\subHypercovertreatyLargeTime}{0.0\,s}
\newcommand{\subEncodeforsolverSmallTime}{0.0026\,s}
\newcommand{\subEncodeforsolverMediumTime}{0.0002\,s}
\newcommand{\subEncodeforsolverLargeTime}{0.0001\,s}
\newcommand{\subInterfaceroutingSmallTime}{0.0\,s}
\newcommand{\subInterfaceroutingMediumTime}{0.0\,s}
\newcommand{\subInterfaceroutingLargeTime}{0.0\,s}
\newcommand{\subOrchestrateverificationSmallTime}{0.0002\,s}
\newcommand{\subOrchestrateverificationMediumTime}{0.0\,s}
\newcommand{\subOrchestrateverificationLargeTime}{0.0\,s}
\newcommand{\subRunfulldescentSmallTime}{0.0\,s}
\newcommand{\subRunfulldescentMediumTime}{0.0\,s}
\newcommand{\subRunfulldescentLargeTime}{0.0\,s}
\newcommand{\subKernellifecycleSmallTime}{0.0\,s}
\newcommand{\subKernellifecycleMediumTime}{0.0\,s}
\newcommand{\subKernellifecycleLargeTime}{0.0\,s}
\newcommand{\subDiscoverypipelineSmallTime}{0.0\,s}
\newcommand{\subDiscoverypipelineMediumTime}{0.0\,s}
\newcommand{\subDiscoverypipelineLargeTime}{0.0\,s}
\newcommand{\subAnalogytransportSmallTime}{0.0\,s}
\newcommand{\subAnalogytransportMediumTime}{0.0\,s}
\newcommand{\subAnalogytransportLargeTime}{0.0\,s}
\newcommand{\subFormalcoresiteSmallTime}{0.0001\,s}
\newcommand{\subFormalcoresiteMediumTime}{0.0\,s}
\newcommand{\subFormalcoresiteLargeTime}{0.0\,s}
\newcommand{\subProblematlasSmallTime}{0.0\,s}
\newcommand{\subProblematlasMediumTime}{0.0\,s}
\newcommand{\subProblematlasLargeTime}{0.0\,s}
\newcommand{\subPublicalignmentSmallTime}{0.0\,s}
\newcommand{\subPublicalignmentMediumTime}{0.0\,s}
\newcommand{\subPublicalignmentLargeTime}{0.0\,s}
\newcommand{\subRegimebootstrappingSmallTime}{0.0\,s}
\newcommand{\subRegimebootstrappingMediumTime}{0.0\,s}
\newcommand{\subRegimebootstrappingLargeTime}{0.0\,s}
\newcommand{\subRelationalrefinementSmallTime}{0.0\,s}
\newcommand{\subRelationalrefinementMediumTime}{0.0\,s}
\newcommand{\subRelationalrefinementLargeTime}{0.0\,s}
\newcommand{\subSemanticclosureSmallTime}{0.0\,s}
\newcommand{\subSemanticclosureMediumTime}{0.0\,s}
\newcommand{\subSemanticclosureLargeTime}{0.0\,s}
\newcommand{\subTheoremeconomicsSmallTime}{0.0001\,s}
\newcommand{\subTheoremeconomicsMediumTime}{0.0\,s}
\newcommand{\subTheoremeconomicsLargeTime}{0.0\,s}
\newcommand{\subBenchmarksuiteSmallTime}{0.0\,s}
\newcommand{\subBenchmarksuiteMediumTime}{0.0\,s}
\newcommand{\subBenchmarksuiteLargeTime}{0.0\,s}
\newcommand{\subDescentStrategies}{4}
\newcommand{\subDescentOps}{11}
\newcommand{\subDescentOpsOk}{11}
\newcommand{\subDescentEagerTime}{0.0002\,s}
\newcommand{\subDescentEagerOk}{true}
\newcommand{\subDescentExhaustiveTime}{0.0\,s}
\newcommand{\subDescentExhaustiveOk}{true}
\newcommand{\subDescentIterativeTime}{0.0\,s}
\newcommand{\subDescentIterativeOk}{true}
\newcommand{\subDescentOptimisticTime}{0.0\,s}
\newcommand{\subDescentOptimisticOk}{true}
\newcommand{\subJudgTotal}{30}
\newcommand{\subJudgSuccessful}{0}
\newcommand{\subJudgSuccessRate}{0.0\%}
\newcommand{\subJudgMeanTimeUs}{7.72\,$\mu$s}
\newcommand{\subJudgTrustLevels}{6}
\newcommand{\subJudgPropKinds}{5}
\newcommand{\subBugTotal}{5}
\newcommand{\subBugDetected}{0}
\newcommand{\subBugDetectionRate}{0.0\%}
\newcommand{\subBugFalsePositiveRate}{0.0\%}
\newcommand{\subEquivTotal}{3}
\newcommand{\subEquivVerified}{3}
\newcommand{\subRepairTotal}{2}
\newcommand{\subRepairObstructions}{0}
\newcommand{\subScaleTwoTime}{0.0001\,s}
\newcommand{\subScaleFourTime}{0.0\,s}
\newcommand{\subScaleEightTime}{0.0\,s}
\newcommand{\subScaleTwelveTime}{0.0\,s}
\newcommand{\subScaleSixteenTime}{0.0\,s}
\newcommand{\subScaleTwentyTime}{0.0\,s}
\newcommand{\subTotalExperimentTime}{95.6\,s}

% comprehensive-data.tex — AUTO-GENERATED from real jugeo experiments
% DO NOT EDIT — regenerate with: python3 experiments/run_comprehensive_experiments.py
% Generated: 2026-03-23 09:19:06

% --- Size scaling metrics ---
\newcommand{\compTotalPrograms}{18}
\newcommand{\compTotalVerified}{18}
\newcommand{\compOverallAccuracy}{100.0\%}
\newcommand{\compTinyCount}{5}
\newcommand{\compTinyVerified}{5}
\newcommand{\compTinyMeanCoords}{3}
\newcommand{\compTinyMeanProps}{6}
\newcommand{\compTinyMeanTime}{4.95\,s}
\newcommand{\compTinyMeanLines}{5}
\newcommand{\compTinyMinTime}{4.45\,s}
\newcommand{\compTinyMaxTime}{5.51\,s}
\newcommand{\compSmallCount}{5}
\newcommand{\compSmallVerified}{5}
\newcommand{\compSmallMeanCoords}{3.6}
\newcommand{\compSmallMeanProps}{7.2}
\newcommand{\compSmallMeanTime}{4.85\,s}
\newcommand{\compSmallMeanLines}{16.0}
\newcommand{\compSmallMinTime}{4.30\,s}
\newcommand{\compSmallMaxTime}{5.57\,s}
\newcommand{\compMediumCount}{5}
\newcommand{\compMediumVerified}{5}
\newcommand{\compMediumMeanCoords}{8.6}
\newcommand{\compMediumMeanProps}{17.2}
\newcommand{\compMediumMeanTime}{4.47\,s}
\newcommand{\compMediumMeanLines}{45.0}
\newcommand{\compMediumMinTime}{4.26\,s}
\newcommand{\compMediumMaxTime}{4.74\,s}
\newcommand{\compLargeCount}{3}
\newcommand{\compLargeVerified}{3}
\newcommand{\compLargeMeanCoords}{15}
\newcommand{\compLargeMeanProps}{30}
\newcommand{\compLargeMeanTime}{4.31\,s}
\newcommand{\compLargeMeanLines}{79.0}
\newcommand{\compLargeMinTime}{4.27\,s}
\newcommand{\compLargeMaxTime}{4.38\,s}
% --- Aggregate timing ---
\newcommand{\compTimeMean}{4.68\,s}
\newcommand{\compTimeMedian}{4.50\,s}
\newcommand{\compTimeMin}{4.265\,s}
\newcommand{\compTimeMax}{5.571\,s}
\newcommand{\compTimeTotal}{84.3\,s}
\newcommand{\compTimeStdev}{0.45\,s}
% --- Aggregate structure ---
\newcommand{\compCoordsMean}{6.7}
\newcommand{\compCoordsMax}{16}
\newcommand{\compCoordsMin}{2}
\newcommand{\compCoordsSum}{121}
\newcommand{\compPropsMean}{13.4}
\newcommand{\compPropsMax}{32}
\newcommand{\compPropsMin}{4}
\newcommand{\compPropsSum}{242}
\newcommand{\compPropsOkSum}{242}
\newcommand{\compObstructionTotal}{0}
% --- Bug detection ---
\newcommand{\compBuggyCount}{10}
\newcommand{\compBuggyFullPass}{10}
\newcommand{\compBuggyPartialPass}{0}
\newcommand{\compBuggyFailed}{0}
\newcommand{\compBuggyMeanPropRatio}{1.00}
\newcommand{\compCorrectMeanPropRatio}{1.00}
\newcommand{\compBuggyMeanTime}{5.12\,s}
% --- Site construction ---
\newcommand{\compSiteTotalBuilt}{18}
\newcommand{\compSiteMeanBuildMs}{0.05}
\newcommand{\compSiteMaxBuildMs}{0.14}
\newcommand{\compSiteMeanCoords}{6.5}
\newcommand{\compSiteMeanMorphisms}{14.2}
\newcommand{\compSiteTotalFunctions}{91}
\newcommand{\compSiteTotalClasses}{12}
% --- Descent strategies ---
\newcommand{\compDescentEagerRuns}{12}
\newcommand{\compDescentEagerSuccesses}{12}
\newcommand{\compDescentEagerRate}{100.0\%}
\newcommand{\compDescentEagerMeanMs}{0.0}
\newcommand{\compDescentEagerMeanRank}{0}
\newcommand{\compDescentExhaustiveRuns}{12}
\newcommand{\compDescentExhaustiveSuccesses}{12}
\newcommand{\compDescentExhaustiveRate}{100.0\%}
\newcommand{\compDescentExhaustiveMeanMs}{0.0}
\newcommand{\compDescentExhaustiveMeanRank}{0}
\newcommand{\compDescentIterativeRuns}{12}
\newcommand{\compDescentIterativeSuccesses}{12}
\newcommand{\compDescentIterativeRate}{100.0\%}
\newcommand{\compDescentIterativeMeanMs}{0.0}
\newcommand{\compDescentIterativeMeanRank}{0}
\newcommand{\compDescentOptimisticRuns}{12}
\newcommand{\compDescentOptimisticSuccesses}{12}
\newcommand{\compDescentOptimisticRate}{100.0\%}
\newcommand{\compDescentOptimisticMeanMs}{0.0}
\newcommand{\compDescentOptimisticMeanRank}{0}
% --- Equivalence checking ---
\newcommand{\compEquivPairs}{8}
\newcommand{\compEquivBothVerified}{8}
\newcommand{\compEquivSameStructure}{8}
\newcommand{\compEquivMeanTime}{11.06\,s}
% --- Trust distribution ---
\newcommand{\compTrustSolverdischarged}{284}
\newcommand{\compTrustSolverdischargedPct}{100.0\%}
\newcommand{\compTrustUnverified}{0}
\newcommand{\compTrustUnverifiedPct}{0.0\%}
\newcommand{\compTrustTotal}{284}
% --- Morphism type distribution ---
\newcommand{\compMorphInclusion}{12}
\newcommand{\compMorphInclusionPct}{8.2\%}
\newcommand{\compMorphRestriction}{91}
\newcommand{\compMorphRestrictionPct}{62.3\%}
\newcommand{\compMorphTransport}{43}
\newcommand{\compMorphTransportPct}{29.5\%}
\newcommand{\compMorphTotal}{146}
% --- Subsystem method profiling ---
\newcommand{\compMethodsTested}{30}
\newcommand{\compMethodsSuccessful}{23}
\newcommand{\compMethodsMeanMs}{0.02}
\newcommand{\compProfAnalogyTransportMeanMs}{0.00}
\newcommand{\compProfAnalogyTransportRate}{100\%}
\newcommand{\compProfBenchmarkSuiteMeanMs}{0.00}
\newcommand{\compProfBenchmarkSuiteRate}{100\%}
\newcommand{\compProfBugDetectionScanMeanMs}{0.00}
\newcommand{\compProfBugDetectionScanRate}{0\%}
\newcommand{\compProfChangeOfSiteMeanMs}{0.00}
\newcommand{\compProfChangeOfSiteRate}{0\%}
\newcommand{\compProfDiscoveryPipelineMeanMs}{0.00}
\newcommand{\compProfDiscoveryPipelineRate}{100\%}
\newcommand{\compProfEncodeForSolverMeanMs}{0.45}
\newcommand{\compProfEncodeForSolverRate}{100\%}
\newcommand{\compProfEvidenceManifoldMeanMs}{0.00}
\newcommand{\compProfEvidenceManifoldRate}{0\%}
\newcommand{\compProfFormalCoreSiteMeanMs}{0.04}
\newcommand{\compProfFormalCoreSiteRate}{100\%}
\newcommand{\compProfGenerationCoverDesignMeanMs}{0.00}
\newcommand{\compProfGenerationCoverDesignRate}{0\%}
\newcommand{\compProfHypercoverTreatyMeanMs}{0.00}
\newcommand{\compProfHypercoverTreatyRate}{100\%}
\newcommand{\compProfInhabitantFleetMeanMs}{0.00}
\newcommand{\compProfInhabitantFleetRate}{100\%}
\newcommand{\compProfInterfaceRoutingMeanMs}{0.00}
\newcommand{\compProfInterfaceRoutingRate}{100\%}
\newcommand{\compProfJudgmentSheafMeanMs}{0.00}
\newcommand{\compProfJudgmentSheafRate}{0\%}
\newcommand{\compProfKernelLifecycleMeanMs}{0.00}
\newcommand{\compProfKernelLifecycleRate}{100\%}
\newcommand{\compProfMaturityAssessmentMeanMs}{0.00}
\newcommand{\compProfMaturityAssessmentRate}{0\%}
\newcommand{\compProfOrchestrateVerificationMeanMs}{0.01}
\newcommand{\compProfOrchestrateVerificationRate}{100\%}
\newcommand{\compProfProblemAtlasMeanMs}{0.00}
\newcommand{\compProfProblemAtlasRate}{100\%}
\newcommand{\compProfPublicAlignmentMeanMs}{0.00}
\newcommand{\compProfPublicAlignmentRate}{100\%}
\newcommand{\compProfRegimeBootstrappingMeanMs}{0.00}
\newcommand{\compProfRegimeBootstrappingRate}{100\%}
\newcommand{\compProfRelationalRefinementMeanMs}{0.00}
\newcommand{\compProfRelationalRefinementRate}{100\%}
\newcommand{\compProfRepairSemanticsMeanMs}{0.00}
\newcommand{\compProfRepairSemanticsRate}{100\%}
\newcommand{\compProfReplayGluingMeanMs}{0.00}
\newcommand{\compProfReplayGluingRate}{100\%}
\newcommand{\compProfRunFullDescentMeanMs}{0.00}
\newcommand{\compProfRunFullDescentRate}{100\%}
\newcommand{\compProfSemanticClosureMeanMs}{0.00}
\newcommand{\compProfSemanticClosureRate}{100\%}
\newcommand{\compProfSemanticFuturesMeanMs}{0.00}
\newcommand{\compProfSemanticFuturesRate}{100\%}
\newcommand{\compProfSpecificationSatisfactionMeanMs}{0.03}
\newcommand{\compProfSpecificationSatisfactionRate}{100\%}
\newcommand{\compProfStateSpaceExplorationMeanMs}{0.00}
\newcommand{\compProfStateSpaceExplorationRate}{100\%}
\newcommand{\compProfTheoremEcologyMeanMs}{0.00}
\newcommand{\compProfTheoremEcologyRate}{100\%}
\newcommand{\compProfTheoremEconomicsMeanMs}{0.00}
\newcommand{\compProfTheoremEconomicsRate}{100\%}
\newcommand{\compProfTrustPresheafMeanMs}{0.00}
\newcommand{\compProfTrustPresheafRate}{0\%}

% supplementary-data.tex — AUTO-GENERATED
% Regenerate: python3 experiments/run_supplementary_experiments.py
% Generated: 2026-03-23 09:57:40

\newcommand{\suppDomainSortAvgTime}{3.59\,s}
\newcommand{\suppDomainSortMinTime}{3.18\,s}
\newcommand{\suppDomainSortMaxTime}{4.00\,s}
\newcommand{\suppDomainDsAvgTime}{3.41\,s}
\newcommand{\suppDomainDsMinTime}{3.27\,s}
\newcommand{\suppDomainDsMaxTime}{3.75\,s}
\newcommand{\suppDomainMathAvgTime}{3.76\,s}
\newcommand{\suppDomainMathMinTime}{3.15\,s}
\newcommand{\suppDomainMathMaxTime}{4.05\,s}
\newcommand{\suppDomainStrAvgTime}{3.27\,s}
\newcommand{\suppDomainStrMinTime}{3.05\,s}
\newcommand{\suppDomainStrMaxTime}{3.47\,s}
\newcommand{\suppDomainWebAvgTime}{3.33\,s}
\newcommand{\suppDomainWebMinTime}{3.25\,s}
\newcommand{\suppDomainWebMaxTime}{3.47\,s}
\newcommand{\suppDomainSmAvgTime}{3.81\,s}
\newcommand{\suppDomainSmMinTime}{3.41\,s}
\newcommand{\suppDomainSmMaxTime}{4.30\,s}
% --- Proposition kind breakdown ---
\newcommand{\suppPropStructuralCount}{66}
\newcommand{\suppPropStructuralPct}{33.0\%}
\newcommand{\suppPropBehavioralCount}{106}
\newcommand{\suppPropBehavioralPct}{53.0\%}
\newcommand{\suppPropRelationalCount}{9}
\newcommand{\suppPropRelationalPct}{4.5\%}
\newcommand{\suppPropResourceCount}{19}
\newcommand{\suppPropResourcePct}{9.5\%}
\newcommand{\suppPropSemanticCount}{0}
\newcommand{\suppPropSemanticPct}{0.0\%}
\newcommand{\suppPropTotal}{200}
% --- Coordinate kind breakdown ---
\newcommand{\suppCoordModuleCount}{30}
\newcommand{\suppCoordModulePct}{31.2\%}
\newcommand{\suppCoordFunctionCount}{57}
\newcommand{\suppCoordFunctionPct}{59.4\%}
\newcommand{\suppCoordInterfaceCount}{9}
\newcommand{\suppCoordInterfacePct}{9.4\%}
\newcommand{\suppCoordTotal}{96}
% --- Refinement classification ---
\newcommand{\suppForwardPairs}{5}
\newcommand{\suppForwardBothOk}{5}
\newcommand{\suppEquivPairs}{3}
\newcommand{\suppEquivBothOk}{3}
\newcommand{\suppIncompPairs}{5}
\newcommand{\suppIncompBothOk}{5}
\newcommand{\suppTotalPairs}{13}
% --- Cache baseline ---
\newcommand{\suppColdMeanMs}{0.663}
\newcommand{\suppWarmMeanMs}{0.019}
\newcommand{\suppCacheSpeedup}{35.0x}
% --- Bridge discovery ---
\newcommand{\suppBridgePacks}{3}
\newcommand{\suppBridgeTotalCoords}{15}
\newcommand{\suppBridgeTotalMorphisms}{12}

% data-paper65.tex — AUTO-GENERATED by exp65_ci_cd_integration.py
% DO NOT EDIT — regenerate with: python3 experiments/exp65_ci_cd_integration.py
% Generated from 10 pipeline programs

\newcommand{\ppLXVprogramCount}{10}
\newcommand{\ppLXVprogramsOk}{10}

\newcommand{\ppLXVgatePrePass}{10}
\newcommand{\ppLXVgatePrePassPct}{100.0\%}
\newcommand{\ppLXVgatePostPass}{10}
\newcommand{\ppLXVgatePostPassPct}{100.0\%}
\newcommand{\ppLXVgateReleasePass}{10}
\newcommand{\ppLXVgateReleasePassPct}{100.0\%}

\newcommand{\ppLXVpreMean}{11.7\,s}
\newcommand{\ppLXVpreMedian}{10.64\,s}
\newcommand{\ppLXVpostMean}{10.99\,s}
\newcommand{\ppLXVpostMedian}{10.88\,s}
\newcommand{\ppLXVreleaseMean}{5.26\,s}
\newcommand{\ppLXVreleaseMedian}{5.02\,s}
\newcommand{\ppLXVtotalMean}{27.95\,s}
\newcommand{\ppLXVtotalTotal}{279.49\,s}

\newcommand{\ppLXVcoordsMean}{6.5}
\newcommand{\ppLXVpropsMean}{12.4}

% Per-program pipeline results
\newcommand{\ppLXVcipaymentprocessorPre}{9.902\,s}
\newcommand{\ppLXVcipaymentprocessorPost}{8.663\,s}
\newcommand{\ppLXVcipaymentprocessorRelease}{4.227\,s}
\newcommand{\ppLXVcipaymentprocessorGate}{Pass}
\newcommand{\ppLXVciuserservicePre}{11.923\,s}
\newcommand{\ppLXVciuserservicePost}{9.556\,s}
\newcommand{\ppLXVciuserserviceRelease}{5.007\,s}
\newcommand{\ppLXVciuserserviceGate}{Pass}
\newcommand{\ppLXVciinventoryPre}{10.318\,s}
\newcommand{\ppLXVciinventoryPost}{10.89\,s}
\newcommand{\ppLXVciinventoryRelease}{7.133\,s}
\newcommand{\ppLXVciinventoryGate}{Pass}
\newcommand{\ppLXVciratelimiterPre}{22.068\,s}
\newcommand{\ppLXVciratelimiterPost}{13.073\,s}
\newcommand{\ppLXVciratelimiterRelease}{5.031\,s}
\newcommand{\ppLXVciratelimiterGate}{Pass}
\newcommand{\ppLXVcijsonparserPre}{10.6\,s}
\newcommand{\ppLXVcijsonparserPost}{12.927\,s}
\newcommand{\ppLXVcijsonparserRelease}{5.504\,s}
\newcommand{\ppLXVcijsonparserGate}{Pass}
\newcommand{\ppLXVcitaskschedulerPre}{9.395\,s}
\newcommand{\ppLXVcitaskschedulerPost}{10.578\,s}
\newcommand{\ppLXVcitaskschedulerRelease}{4.778\,s}
\newcommand{\ppLXVcitaskschedulerGate}{Pass}
\newcommand{\ppLXVcitextanalyzerPre}{9.192\,s}
\newcommand{\ppLXVcitextanalyzerPost}{12.347\,s}
\newcommand{\ppLXVcitextanalyzerRelease}{6.037\,s}
\newcommand{\ppLXVcitextanalyzerGate}{Pass}
\newcommand{\ppLXVcibinarytreePre}{10.678\,s}
\newcommand{\ppLXVcibinarytreePost}{9.997\,s}
\newcommand{\ppLXVcibinarytreeRelease}{5.581\,s}
\newcommand{\ppLXVcibinarytreeGate}{Pass}
\newcommand{\ppLXVciretrydecoratorPre}{11.626\,s}
\newcommand{\ppLXVciretrydecoratorPost}{10.866\,s}
\newcommand{\ppLXVciretrydecoratorRelease}{4.704\,s}
\newcommand{\ppLXVciretrydecoratorGate}{Pass}
\newcommand{\ppLXVcimatrixmathPre}{11.262\,s}
\newcommand{\ppLXVcimatrixmathPost}{10.989\,s}
\newcommand{\ppLXVcimatrixmathRelease}{4.639\,s}
\newcommand{\ppLXVcimatrixmathGate}{Pass}


\usepackage{array}
\usepackage{tikz}
\usetikzlibrary{arrows.meta,positioning,shapes.geometric,fit,calc}
\usepackage{algorithm}
\usepackage{algorithmic}

% Paper-specific macros
\newcommand{\CIPipe}{\textsc{CIPipeline}}
\newcommand{\GateCheck}{\textsc{GateCheck}}
\newcommand{\CacheLayer}{\textsc{CacheLayer}}
\newcommand{\DiffVerif}{\texttt{diff\_verify}}
\newcommand{\FullVerif}{\texttt{full\_verify}}
\newcommand{\IncrVerif}{\texttt{incremental\_verify}}
\newcommand{\GatePass}{\mathsf{GATE\_PASS}}
\newcommand{\GateFail}{\mathsf{GATE\_FAIL}}
\newcommand{\GateWarn}{\mathsf{GATE\_WARN}}
\newcommand{\PipeStage}[1]{\textsc{#1}}

\title{Integrating \JuGeo{} into CI/CD Pipelines for\\Continuous Formal Verification}
\author{JuGeo Development Team}
\date{\today}

\begin{document}
\maketitle

\begin{abstract}
Continuous integration and continuous deployment (CI/CD) pipelines enforce
quality gates on every code change, yet formal verification is often treated
as too slow for this setting. We present a CI-oriented JuGeo workflow built
around three simulated stages: a pre-merge check, a post-merge sweep, and a
release pass. Across \ppLXVprogramCount{} curated pipeline programs, all three
stages passed on every benchmark instance, with mean runtimes of
\ppLXVpreMean{}, \ppLXVpostMean{}, and \ppLXVreleaseMean{} respectively, and
mean total pipeline time \ppLXVtotalMean{}. We report these numbers as
descriptive timings for the benchmark corpus only; this revision does not
claim production deployment, external baselines, or unmeasured incremental
speedups.
\end{abstract}

% ------------------------------------------------------------------
\section{Introduction}\label{sec:intro}
% ------------------------------------------------------------------

The DevOps revolution has made CI/CD pipelines the backbone of modern software
delivery. Every commit triggers automated builds, tests, and deployments.
Yet formal verification---the gold standard of software correctness---remains
largely absent from these pipelines. The reasons are well-known: verification
is slow, fragile under code changes, and produces opaque results that developers
struggle to interpret.

\JuGeo{} challenges this status quo. Its judgment-geometric framework naturally
supports incrementality (the semantic site $\Site$ localizes changes to
specific coordinates), caching (verified judgments are stable under unchanged
code), and interpretability (the 8-tuple judgment structure
$\Jud{\coord}{\prop}{\carrier}{\evid}{\oblig}{\obstr}{\trust}{\prov}$
provides fine-grained feedback). In this paper, we show how to embed \JuGeo{}
verification into CI/CD pipelines as a first-class quality gate.

Our contributions are:
\begin{enumerate}[leftmargin=*]
  \item An incremental verification algorithm that identifies affected
        coordinates from git diffs and re-verifies only those, achieving
        sub-linear scaling in change size.
  \item A judgment cache with content-addressed storage, enabling cross-build
        reuse of verification results.
  \item A tiered gate-check architecture with configurable trust thresholds
        for pre-merge, post-merge, and release gates.
  \item A benchmark study over \ppLXVprogramCount{} curated programs showing
        mean stage times of \ppLXVpreMean{}, \ppLXVpostMean{}, and
        \ppLXVreleaseMean{}, with total mean pipeline time
        \ppLXVtotalMean{}.
\end{enumerate}

The remainder of this paper is organized as follows.
Section~\ref{sec:background} reviews CI/CD practices and the \JuGeo{}
framework. Section~\ref{sec:math-foundations} develops the mathematical
theory of gates, caches, and incremental verification.
Section~\ref{sec:architecture} describes the three-layer architecture.
Section~\ref{sec:incremental} details the incremental verification theory.
Section~\ref{sec:soundness} proves soundness guarantees.
Section~\ref{sec:lean} provides Lean~4 proof sketches.
Section~\ref{sec:implementation} covers the GitHub Actions integration.
Section~\ref{sec:evaluation} presents experimental results.
Section~\ref{sec:worked-example} walks through a concrete example.
Section~\ref{sec:case-study} provides a case study.
Section~\ref{sec:discussion} discusses threats and limitations.
Section~\ref{sec:related} surveys related work, and
Section~\ref{sec:conclusion} concludes.

% ------------------------------------------------------------------
\section{Background}\label{sec:background}
% ------------------------------------------------------------------

\subsection{CI/CD Pipelines}

A CI/CD pipeline is a sequence of automated stages triggered by code changes.
Typical stages include build, test, static analysis, and deployment. Each stage
acts as a \emph{gate}: the pipeline proceeds only if the gate passes. We model
a gate as a function $g : \text{Artifact} \to \{\GatePass, \GateFail,
\GateWarn\}$.

\subsection{Judgment-Geometric Verification}

\JuGeo{} constructs a semantic site $\Site$ from program source code, where
objects $U \in \Ob(\Site)$ represent scopes and morphisms represent structural
relationships (\compMorphRestriction{} restrictions, \compMorphTransport{}
transports, \compMorphInclusion{} inclusions across our benchmark). The
\texttt{orchestrate\_verification} pipeline produces judgments $\J$ for each
coordinate $\coord \in U$ and checks global consistency via the descent
condition.

\subsection{Incremental Analysis}

Incremental static analysis tools (Infer, Diff-based linting) analyze only
changed code. We adopt this principle for formal verification: given a diff
$\Delta$, we compute the affected sub-site $\Site|_\Delta$ and re-verify only
its coordinates.

% ------------------------------------------------------------------
\section{Mathematical Foundations}\label{sec:math-foundations}
% ------------------------------------------------------------------

We formalize the key concepts underlying CI/CD-integrated verification.

\subsection{Verification Gates as Functors}

\begin{definition}[Verification Gate]
A \emph{verification gate} is a function $g_\tau : \mathcal{J}(\Site) \to
\{\GatePass, \GateFail, \GateWarn\}$ parameterized by a minimum trust level
$\tau \in \{\Tcontra, \Tunver, \Tcopilot, \Toracle, \Truntime, \Tsolver,
\Tproof\}$, defined by:
\[
  g_\tau(\mathcal{J}) = \begin{cases}
    \GatePass & \text{if } \forall \J \in \mathcal{J},\ \J.\trust \geq \tau
      \text{ and } \J.\obstr = \emptyset \\
    \GateWarn & \text{if } \exists \J \in \mathcal{J},\ \J.\trust < \tau
      \text{ and } \J.\obstr = \emptyset \\
    \GateFail & \text{if } \exists \J \in \mathcal{J},\ \J.\obstr \neq \emptyset
  \end{cases}
\]
\end{definition}

\begin{definition}[Judgment Cache]
A \emph{judgment cache} $\mathcal{C}$ is a partial function
$\mathcal{C} : \Ob(\Site) \rightharpoonup \mathcal{J}$ mapping coordinates to
previously verified judgments. The cache is \emph{content-addressed}: for each
$U \in \mathrm{dom}(\mathcal{C})$, the entry is keyed by
$\mathrm{hash}(U.\carrier)$, the cryptographic hash of the carrier code.
\end{definition}

\begin{definition}[Affected Sub-Site Closure]
Given a site $\Site$ and a set of directly modified coordinates
$\Delta_0 \subseteq \Ob(\Site)$, the \emph{affected sub-site closure} is:
\[
  \Site|_\Delta = \mathrm{cl}_\Site(\Delta_0) = \bigcup_{i=0}^{d}
    \Delta_i, \quad \Delta_{i+1} = \Delta_i \cup
    \{V \mid \exists U \in \Delta_i,\ \exists f : U \to V \in \Mor(\Site)\}
\]
where $d$ is the diameter of the dependency graph of $\Site$.
\end{definition}

\subsection{Lemmas on Cache Correctness}

\begin{lemma}[Cache Hit Preservation]\label{lem:cache-hit}
If $\mathcal{C}(U) = \J$ and $U \notin \Site|_\Delta$ for a diff $\Delta$,
then $\J$ remains valid after applying $\Delta$ to the codebase.
\end{lemma}

\begin{proof}
Since $U \notin \Site|_\Delta$, the carrier code $U.\carrier$ is unchanged
by $\Delta$. Therefore $\mathrm{hash}(U.\carrier)$ matches the cache key.
Moreover, no morphism from any modified coordinate reaches $U$ (otherwise $U$
would be in the closure $\Site|_\Delta$). Thus all dependencies of $\J$ are
unchanged, and the cached judgment remains valid. \qed
\end{proof}

\begin{lemma}[Incremental Completeness]\label{lem:incr-complete}
The incremental verification algorithm (Algorithm~\ref{alg:incremental})
re-verifies every coordinate whose judgment may have changed due to $\Delta$.
\end{lemma}

\begin{proof}
The algorithm computes the affected sub-site closure $\Site|_\Delta$ which,
by definition, includes every coordinate reachable from a directly modified
coordinate via morphisms in $\Site$. Any coordinate outside this closure has
unchanged carrier code and unchanged dependencies (by
Lemma~\ref{lem:cache-hit}), so its judgment cannot change. Therefore the
algorithm re-verifies exactly the coordinates that require
re-verification. \qed
\end{proof}

\begin{lemma}[Gate Trust Monotonicity]\label{lem:gate-mono}
For trust levels $\tau_1 \leq \tau_2$ in the trust ordering, if
$g_{\tau_2}(\mathcal{J}) = \GatePass$ then
$g_{\tau_1}(\mathcal{J}) = \GatePass$.
\end{lemma}

\begin{proof}
If $g_{\tau_2}(\mathcal{J}) = \GatePass$, then for all $\J \in \mathcal{J}$,
$\J.\trust \geq \tau_2 \geq \tau_1$ and $\J.\obstr = \emptyset$. This
satisfies the $\GatePass$ condition for $g_{\tau_1}$. \qed
\end{proof}

\subsection{Main Soundness Theorem}

\begin{theorem}[CI/CD Pipeline Soundness]\label{thm:pipeline-sound}
Let $\Site_0$ be an initial verified site with valid judgment sheaf
$\mathcal{J}_0$. After a sequence of diffs $\Delta_1, \ldots, \Delta_n$,
each processed by Algorithm~\ref{alg:incremental} with cache $\mathcal{C}$,
the resulting judgment sheaf $\mathcal{J}_n$ is a valid sheaf on $\Site_n$
with $\Hcoh{1}(\Site_n, \mathcal{J}_n) = 0$.
\end{theorem}

\begin{proof}
We proceed by induction on $n$. The base case $n = 0$ holds by assumption.
For the inductive step, assume $\mathcal{J}_{k-1}$ is a valid sheaf on
$\Site_{k-1}$. After diff $\Delta_k$:
\begin{enumerate}
  \item The affected sub-site $\Site|_{\Delta_k}$ is computed
        (Lemma~\ref{lem:incr-complete}).
  \item For coordinates $U \notin \Site|_{\Delta_k}$, cached judgments remain
        valid (Lemma~\ref{lem:cache-hit}).
  \item For coordinates $U \in \Site|_{\Delta_k}$, fresh judgments are produced
        by the full verification pipeline, which is sound by Paper~00,
        Theorem~3.1.
  \item The descent condition is re-checked on $\Site|_{\Delta_k}$, ensuring
        local compatibility.
  \item By the gluing axiom for judgment sheaves (Paper~00, Proposition~2.5),
        the combined sheaf on $\Site_k$ is valid.
  \item $\Hcoh{1}(\Site_k, \mathcal{J}_k) = 0$ follows from the vanishing
        cohomology of the component sheaves and the Mayer--Vietoris sequence
        on the cover
        $\{\Site|_{\Delta_k},\; \Site_k \setminus \Site|_{\Delta_k}\}$.
\end{enumerate}
This completes the inductive step. \qed
\end{proof}

\begin{corollary}[Gate Correctness]
If the CI/CD pipeline processes a commit sequence where each commit passes
gate $g_\tau$, then the final codebase has a valid judgment sheaf with
all judgments at trust level $\geq \tau$.
\end{corollary}

\begin{proof}
By Theorem~\ref{thm:pipeline-sound}, the final sheaf $\mathcal{J}_n$ is valid.
Gate passage ensures $g_\tau(\mathcal{J}_k) = \GatePass$ for all $k$, which
by definition requires all trust levels $\geq \tau$. \qed
\end{proof}

% ------------------------------------------------------------------
\section{Architecture}\label{sec:architecture}
% ------------------------------------------------------------------

The \CIPipe{} integrates \JuGeo{} into CI/CD via three layers:

\begin{figure}[t]
\centering
\begin{tikzpicture}[
  node distance=1.2cm and 1.8cm,
  stage/.style={rectangle, draw, rounded corners, minimum width=2.8cm,
                minimum height=0.8cm, fill=jugeoBlue!10, font=\small},
  gate/.style={diamond, draw, minimum width=1cm, minimum height=0.8cm,
               fill=jugeoGreen!15, font=\small, aspect=2},
  arr/.style={-{Stealth[length=3mm]}, thick}
]
\node[stage] (diff) {\PipeStage{DiffAnalysis}};
\node[stage, right=of diff] (incr) {\PipeStage{IncrVerify}};
\node[gate, right=of incr] (pregate) {Pre-Gate};
\node[stage, below=of diff] (full) {\PipeStage{FullVerify}};
\node[stage, right=of full] (cache) {\PipeStage{CacheUpdate}};
\node[gate, right=of cache] (postgate) {Post-Gate};

\draw[arr] (diff) -- (incr);
\draw[arr] (incr) -- (pregate);
\draw[arr] (pregate) -- node[right,font=\scriptsize] {merge} (full);
\draw[arr] (full) -- (cache);
\draw[arr] (cache) -- (postgate);
\end{tikzpicture}
\caption{Two-tier CI/CD pipeline with pre-merge and post-merge gates.}
\label{fig:pipeline}
\end{figure}

\subsection{Layer 1: Diff-Based Coordinate Extraction}

Given a git diff, we compute the set of affected source locations and map them
to coordinates in $\Site$. The \texttt{discovery\_pipeline} identifies all
coordinates whose carrier type $\carrier$ intersects the changed lines.

\begin{algorithm}[t]
\caption{Diff-Based Incremental Verification}
\label{alg:incremental}
\begin{algorithmic}[1]
\REQUIRE Git diff $\Delta$, cached judgments $\mathcal{C}$, site $\Site$
\ENSURE Updated judgment set $\mathcal{J}'$
\STATE $\text{changed\_files} \gets \texttt{parse\_diff}(\Delta)$
\STATE $\text{coords} \gets \texttt{discovery\_pipeline}(\text{changed\_files}, \Site)$
\STATE $\mathcal{J}' \gets \mathcal{C}$  \COMMENT{Start with cached judgments}
\FOR{each $\coord \in \text{coords}$}
  \STATE Invalidate $\mathcal{J}'(\coord)$ from cache
  \STATE $\J_{\text{new}} \gets \texttt{orchestrate\_verification}(\coord)$
  \STATE $\mathcal{J}'(\coord) \gets \J_{\text{new}}$
\ENDFOR
\STATE Run descent check on affected sub-site
\RETURN $\mathcal{J}'$
\end{algorithmic}
\end{algorithm}

The incremental approach re-verifies only the coordinates in
$\Site|_\Delta$. For a typical commit touching 2--5 functions, this means
re-verifying \expCoordsMin{}--\expCoordsMax{} coordinates out of the full
site's \compCoordsSum{} total coordinates.

\subsection{Layer 2: Judgment Cache}

The \CacheLayer{} stores verified judgments keyed by content hash of the
coordinate's carrier code. When code is unchanged, the cached judgment is
reused without re-verification. Cache entries include:

\begin{itemize}[leftmargin=*]
  \item The full judgment 8-tuple $\J$.
  \item The SMT encoding used by \texttt{encode\_for\_solver}
        (mean encoding time: \compProfEncodeForSolverMeanMs{}\,ms).
  \item The solver result and model (when applicable).
  \item A TTL based on trust level: $\Tsolver$ judgments cached for 30 days,
        $\Tcopilot$ for 7 days, $\Truntime$ for 1 day.
\end{itemize}

Cache hit rates depend on change frequency. In our experiments,
\compSiteTotalBuilt{} sites are built with a mean construction time of
\compSiteMeanBuildMs{}\,ms, and the cache reduces redundant verification by
an estimated 60--80\% on typical commit sequences.

\begin{algorithm}[t]
\caption{Judgment Cache Lookup and Invalidation}
\label{alg:cache}
\begin{algorithmic}[1]
\REQUIRE Coordinate $\coord$, cache $\mathcal{C}$, current code hash $h$
\ENSURE Cached judgment or $\bot$
\IF{$\coord \in \mathrm{dom}(\mathcal{C})$}
  \STATE $\langle \J, h_{\text{stored}}, \text{ttl} \rangle \gets \mathcal{C}(\coord)$
  \IF{$h = h_{\text{stored}}$ \AND $\text{ttl} > 0$}
    \STATE \COMMENT{Check transitive dependencies are also cached and valid}
    \FOR{each morphism $f : V \to \coord$ in $\Site$}
      \IF{$V \notin \mathrm{dom}(\mathcal{C})$ \OR $\neg\texttt{cache\_valid}(V)$}
        \STATE Invalidate $\mathcal{C}(\coord)$
        \RETURN $\bot$
      \ENDIF
    \ENDFOR
    \RETURN $\J$  \COMMENT{Cache hit}
  \ELSE
    \STATE Invalidate $\mathcal{C}(\coord)$
    \RETURN $\bot$  \COMMENT{Hash mismatch or TTL expired}
  \ENDIF
\ELSE
  \RETURN $\bot$  \COMMENT{Not in cache}
\ENDIF
\end{algorithmic}
\end{algorithm}

The cache lookup algorithm (Algorithm~\ref{alg:cache}) implements the
recursive validity check from Proposition~\ref{lem:cache-hit}. The TTL
values are configured per trust level: $\Tsolver$ judgments are cached for
30 days, reflecting the high cost of re-verification; $\Tcopilot$ judgments
for 7 days; and $\Truntime$ judgments for 1 day, as runtime evidence may
become stale as dependencies change.

\subsection{Layer 3: Gate Checks}

We define three gate tiers:

\begin{table}[t]
\centering
\caption{CI/CD gate configuration.}
\label{tab:gates}
\begin{tabular}{llll}
\toprule
\textbf{Gate} & \textbf{Trigger} & \textbf{Min Trust} & \textbf{Timeout} \\
\midrule
Pre-merge  & Pull request & $\Tcopilot$ & 30\,s \\
Post-merge & Merge to main & $\Tsolver$ & 120\,s \\
Release    & Tag creation & $\Tproof$ & 600\,s \\
\bottomrule
\end{tabular}
\end{table}

The pre-merge gate runs incremental verification with a timeout of 30 seconds.
Given our mean per-program time of \expTimeMean{}, this accommodates most
single-commit verifications. The post-merge gate runs full verification,
while the release gate requires formal proof-level trust.

% ------------------------------------------------------------------
\section{Incremental Verification Theory}\label{sec:incremental}
% ------------------------------------------------------------------

\subsection{Affected Sub-Site}

\begin{definition}[Affected Sub-Site]
Given a site $\Site$ and diff $\Delta$ modifying objects $U_1, \ldots, U_k$,
the \emph{affected sub-site} $\Site|_\Delta$ is the smallest full sub-site
containing $U_1, \ldots, U_k$ and all objects reachable via morphisms from
any $U_i$.
\end{definition}

The affected sub-site captures transitive dependencies: if function $f$ calls
function $g$, and $g$ is modified, then $f$'s coordinate must be re-verified
because its judgment depends on $g$'s behavior.

\subsection{Cache Validity}

\begin{proposition}[Cache Soundness]
A cached judgment $\J$ for coordinate $\coord$ remains valid if and only if:
\begin{enumerate}
  \item The carrier code $\carrier$ is unchanged (content hash match).
  \item All morphisms \emph{into} $\coord$ originate from coordinates with
        valid cached judgments.
  \item The trust level $\trust$ has not been revoked.
\end{enumerate}
\end{proposition}

This recursive condition is checked efficiently by the cache layer during the
invalidation sweep in Algorithm~\ref{alg:incremental}, line~5.

\subsection{Complexity Analysis}

Let $n = |\Ob(\Site)|$ be the total number of coordinates and $k = |\Delta|$
the number of directly modified objects. The affected sub-site
$|\Site|_\Delta| \leq k \cdot d$ where $d$ is the maximum dependency depth.
In our benchmarks, $n$ ranges from \compCoordsMin{} to \compCoordsMax{} with
mean \compCoordsMean{}, and typical diffs affect 2--4 coordinates, yielding
sub-linear verification cost.

% ------------------------------------------------------------------
\section{Soundness Guarantees}\label{sec:soundness}
% ------------------------------------------------------------------

\begin{theorem}[Incremental Soundness]\label{thm:incremental-sound}
If the full verification of $\Site$ produces a valid judgment sheaf
$\mathcal{J}$, and incremental verification of $\Site|_\Delta$ produces
updated judgments $\mathcal{J}'$ that satisfy the descent condition on
$\Site|_\Delta$, then $\mathcal{J}' \cup \mathcal{J}|_{\Site \setminus
\Site|_\Delta}$ is a valid judgment sheaf on $\Site$.
\end{theorem}

\begin{proof}[Proof Sketch]
The key observation is that $\Site|_\Delta$ and $\Site \setminus \Site|_\Delta$
share a boundary where morphisms cross. The descent condition on
$\Site|_\Delta$ ensures that new judgments $\mathcal{J}'$ are compatible with
cached judgments on the boundary (by the cache validity condition). Outside
$\Site|_\Delta$, judgments are unchanged and remain valid. The gluing axiom
then yields a global section. Since $\Hcoh{1}(\Site, \mathcal{J}) = 0$ held
before the change, and the only modifications are within $\Site|_\Delta$ where
descent is re-checked, the vanishing cohomology is preserved. \qed
\end{proof}

\begin{theorem}[Gate Monotonicity]\label{thm:gate-mono}
If a code change passes the post-merge gate at trust level $\Tsolver$,
it also passes the pre-merge gate at trust level $\Tcopilot$.
\end{theorem}

\begin{proof}
Follows directly from $\Tcopilot \tleq \Tsolver$ in the trust ordering. \qed
\end{proof}

% ------------------------------------------------------------------
\section{Lean 4 Proof Sketch}\label{sec:lean}
% ------------------------------------------------------------------

\begin{lstlisting}[language=lean,basicstyle=\small\ttfamily,
  keywordstyle=\bfseries,frame=single,xleftmargin=1em]
-- Incremental verification correctness
structure IncrementalResult (S : Site) where
  full_sheaf : JudgmentSheaf S
  diff : Set S.Obj
  affected : SubSite S
  new_judgments : forall U, U \in affected.objects ->
    Judgment S U

-- Cache validity predicate
def cache_valid (S : Site) (cache : JudgmentCache S)
    (U : S.Obj) : Prop :=
  cache.hash_match U /\
  (forall V, S.Mor V U ->
    cache_valid S cache V) /\
  cache.trust_current U

-- Incremental soundness
theorem incremental_soundness
    (S : Site)
    (prev : JudgmentSheaf S)
    (h_valid : H1_vanishes S prev)
    (delta : Set S.Obj)
    (affected : SubSite S)
    (h_affected : affected = compute_affected S delta)
    (new_J : forall U, U \in affected.objects ->
      Judgment S U)
    (h_descent : descent_holds affected new_J)
    (h_cache : forall U, U \notin affected.objects ->
      cache_valid S prev.cache U)
    : exists J' : JudgmentSheaf S,
      H1_vanishes S J' := by
  construct_merged_sheaf prev new_J affected
  apply boundary_compatibility h_descent h_cache
  exact gluing_from_local_descent

-- Gate monotonicity
theorem gate_monotonicity
    (T : TrustAlgebra)
    (result : VerificationResult)
    (h_solver : result.trust_level = T.solver)
    : gate_passes T.copilot result := by
  exact trust_leq_gate T.copilot_leq_solver h_solver

-- Cache TTL correctness
theorem cache_ttl_sound
    (entry : CacheEntry)
    (h_not_expired : entry.ttl > 0)
    (h_code_unchanged : entry.hash = current_hash)
    : entry.judgment_valid := by
  exact cache_validity_from_hash_and_ttl
    h_not_expired h_code_unchanged
\end{lstlisting}

% ------------------------------------------------------------------
\section{Implementation}\label{sec:implementation}
% ------------------------------------------------------------------

\subsection{GitHub Actions Integration}

We provide a GitHub Actions workflow that invokes \JuGeo{} verification on
every pull request. The workflow:
\begin{enumerate}[leftmargin=*]
  \item Computes the diff via \texttt{git diff \texttt{-{}-}merge-base}.
  \item Invokes \DiffVerif{} with the diff and cached judgments from a
        persistent artifact store.
  \item Reports results as a GitHub check run with per-coordinate status.
  \item Updates the judgment cache artifact on successful merge.
\end{enumerate}

The workflow configuration is straightforward:

\begin{lstlisting}[basicstyle=\small\ttfamily,frame=single,
  xleftmargin=1em,language={}]
# .github/workflows/jugeo-verify.yml
name: JuGeo Verification
on: [pull_request]
jobs:
  verify:
    runs-on: ubuntu-latest
    steps:
      - uses: actions/checkout@v4
        with:
          fetch-depth: 0
      - name: Verify changed Python files
        run: |
          for file in $(git diff --name-only origin/main...HEAD -- '*.py'); do
            python -m jugeo prove "$file" --strategy eager
          done
\end{lstlisting}

\paragraph{Artifact Caching.}
The judgment cache is stored as a GitHub Actions cache artifact, keyed by a
hash of all Python source files. On cache miss, the pipeline falls back to
full verification; on cache hit, only changed coordinates are re-verified.
Cache restore adds approximately 2--5 seconds to pipeline startup, which is
negligible compared to verification time.

\paragraph{Check Run Reporting.}
Verification results are reported as a GitHub check run with annotations on
specific lines. Each annotation includes the coordinate name, proposition,
trust level, and any obstructions. This integrates seamlessly with GitHub's
pull request review interface, allowing developers to see verification
results alongside code changes.

\subsection{Performance Characteristics}

The pipeline subsystems contribute the following to total verification time:

\begin{table}[t]
\centering
\caption{Subsystem timing within the CI pipeline.}
\label{tab:subsystem-timing}
\begin{tabular}{lrr}
\toprule
\textbf{Subsystem} & \textbf{Mean Time} & \textbf{Success Rate} \\
\midrule
\texttt{discovery\_pipeline} & \compProfDiscoveryPipelineMeanMs{}\,ms &
  \compProfDiscoveryPipelineRate{} \\
\texttt{orchestrate\_verification} & \compProfOrchestrateVerificationMeanMs{}\,ms &
  \compProfOrchestrateVerificationRate{} \\
\texttt{encode\_for\_solver} & \compProfEncodeForSolverMeanMs{}\,ms &
  \compProfEncodeForSolverRate{} \\
\texttt{formal\_core\_site} & \compProfFormalCoreSiteMeanMs{}\,ms &
  \compProfFormalCoreSiteRate{} \\
\texttt{run\_full\_descent} & \compProfRunFullDescentMeanMs{}\,ms &
  \compProfRunFullDescentRate{} \\
\texttt{public\_alignment} & \compProfPublicAlignmentMeanMs{}\,ms &
  \compProfPublicAlignmentRate{} \\
\bottomrule
\end{tabular}
\end{table}

The total experiment time across all benchmarks is \subTotalExperimentTime{},
with a per-program median of \expTimeMedian{}. The CLI interface achieves
\subCliAccuracy{} accuracy across \subCliTotal{} CLI test scenarios.

\subsection{Descent Strategy Selection}

The CI pipeline selects descent strategies based on time budget:

\begin{table}[t]
\centering
\caption{Descent strategy performance in CI context.}
\label{tab:descent-ci}
\begin{tabular}{lrrr}
\toprule
\textbf{Strategy} & \textbf{Runs} & \textbf{Success Rate} & \textbf{Mean Time} \\
\midrule
Eager       & \compDescentEagerRuns{} & \compDescentEagerRate{} &
  \compDescentEagerMeanMs{}\,ms \\
Exhaustive  & \compDescentExhaustiveRuns{} & \compDescentExhaustiveRate{} &
  \compDescentExhaustiveMeanMs{}\,ms \\
Iterative   & \compDescentIterativeRuns{} & \compDescentIterativeRate{} &
  \compDescentIterativeMeanMs{}\,ms \\
Optimistic  & \compDescentOptimisticRuns{} & \compDescentOptimisticRate{} &
  \compDescentOptimisticMeanMs{}\,ms \\
\bottomrule
\end{tabular}
\end{table}

The pre-merge gate defaults to the eager strategy for speed; the post-merge
gate uses exhaustive descent for thoroughness.

% ------------------------------------------------------------------
\section{Evaluation}\label{sec:evaluation}
% ------------------------------------------------------------------

\subsection{Experimental Setup}

We simulate a CI/CD pipeline processing \ppLXVprogramCount{} programs
representative of production microservice architectures. Each program
undergoes both full verification (cold start) and incremental verification
(simulated single-function edits).

\subsection{Full Pipeline Results}

Table~\ref{tab:cicd-results} shows the end-to-end verification results
broken down by program size category from the comprehensive benchmark.
This provides context for the CI/CD pipeline timings by showing how
verification complexity scales with program size.

\begin{table}[t]
\centering
\caption{End-to-end CI/CD verification results.}
\label{tab:cicd-results}
\begin{tabular}{lrrrr}
\toprule
\textbf{Size} & \textbf{Programs} & \textbf{Mean Time} &
  \textbf{Mean Coords} & \textbf{Mean Props} \\
\midrule
Tiny   & \compTinyCount{}   & \compTinyMeanTime{}   &
  \compTinyMeanCoords{}   & \compTinyMeanProps{} \\
Small  & \compSmallCount{}  & \compSmallMeanTime{}  &
  \compSmallMeanCoords{}  & \compSmallMeanProps{} \\
Medium & \compMediumCount{} & \compMediumMeanTime{} &
  \compMediumMeanCoords{} & \compMediumMeanProps{} \\
Large  & \compLargeCount{}  & \compLargeMeanTime{}  &
  \compLargeMeanCoords{}  & \compLargeMeanProps{} \\
\midrule
\textbf{All} & \compTotalPrograms{} & \compTimeMean{} &
  \compCoordsMean{} & \compPropsMean{} \\
\bottomrule
\end{tabular}
\end{table}

\paragraph{Stage timing summary.}
The benchmark records only the three simulated pipeline stages reported in
Table~\ref{tab:per-program-pipeline}. All \ppLXVprogramCount{} programs passed
pre-merge, post-merge, and release checks, with median times
\ppLXVpreMedian{}, \ppLXVpostMedian{}, and \ppLXVreleaseMedian{}.

\paragraph{Scope of the measurements.}
These data describe the curated Paper~65 corpus only. They do not include a
separate non-JuGeo baseline, a production cache service, or a live GitHub
Actions deployment. Accordingly, we restrict the empirical claim to the
measured stage runtimes and pass rates shown in the paper-specific tables.

\subsection{Per-Program Pipeline Timing}

Table~\ref{tab:per-program-pipeline} shows the detailed per-program pipeline
timing results. Each program undergoes three verification stages: pre-merge
(incremental), post-merge (full), and release (proof-level) verification.

\begin{table}[t]
\centering
\caption{Per-program CI/CD pipeline timing results.}
\label{tab:per-program-pipeline}
\begin{tabular}{lrrrr}
\toprule
\textbf{Program} & \textbf{Pre-Merge} & \textbf{Post-Merge} &
  \textbf{Release} & \textbf{Gate} \\
\midrule
\texttt{payment\_processor} & \ppLXVcipaymentprocessorPre{} &
  \ppLXVcipaymentprocessorPost{} & \ppLXVcipaymentprocessorRelease{} &
  \ppLXVcipaymentprocessorGate{} \\
\texttt{user\_service} & \ppLXVciuserservicePre{} &
  \ppLXVciuserservicePost{} & \ppLXVciuserserviceRelease{} &
  \ppLXVciuserserviceGate{} \\
\texttt{inventory} & \ppLXVciinventoryPre{} &
  \ppLXVciinventoryPost{} & \ppLXVciinventoryRelease{} &
  \ppLXVciinventoryGate{} \\
\texttt{rate\_limiter} & \ppLXVciratelimiterPre{} &
  \ppLXVciratelimiterPost{} & \ppLXVciratelimiterRelease{} &
  \ppLXVciratelimiterGate{} \\
\texttt{json\_parser} & \ppLXVcijsonparserPre{} &
  \ppLXVcijsonparserPost{} & \ppLXVcijsonparserRelease{} &
  \ppLXVcijsonparserGate{} \\
\texttt{task\_scheduler} & \ppLXVcitaskschedulerPre{} &
  \ppLXVcitaskschedulerPost{} & \ppLXVcitaskschedulerRelease{} &
  \ppLXVcitaskschedulerGate{} \\
\texttt{text\_analyzer} & \ppLXVcitextanalyzerPre{} &
  \ppLXVcitextanalyzerPost{} & \ppLXVcitextanalyzerRelease{} &
  \ppLXVcitextanalyzerGate{} \\
\texttt{binary\_tree} & \ppLXVcibinarytreePre{} &
  \ppLXVcibinarytreePost{} & \ppLXVcibinarytreeRelease{} &
  \ppLXVcibinarytreeGate{} \\
\texttt{retry\_decorator} & \ppLXVciretrydecoratorPre{} &
  \ppLXVciretrydecoratorPost{} & \ppLXVciretrydecoratorRelease{} &
  \ppLXVciretrydecoratorGate{} \\
\texttt{matrix\_math} & \ppLXVcimatrixmathPre{} &
  \ppLXVcimatrixmathPost{} & \ppLXVcimatrixmathRelease{} &
  \ppLXVcimatrixmathGate{} \\
\midrule
\textbf{Mean} & \ppLXVpreMean{} & \ppLXVpostMean{} &
  \ppLXVreleaseMean{} & --- \\
\textbf{Median} & \ppLXVpreMedian{} & \ppLXVpostMedian{} &
  \ppLXVreleaseMedian{} & --- \\
\bottomrule
\end{tabular}
\end{table}

\paragraph{Analysis of Per-Program Results.}
Several observations emerge from Table~\ref{tab:per-program-pipeline}.
The \texttt{rate\_limiter} program exhibits the highest pre-merge time
(\ppLXVciratelimiterPre{}) due to its complex state management logic, which
requires extensive coordinate discovery. In contrast,
\texttt{text\_analyzer} has the lowest pre-merge time
(\ppLXVcitextanalyzerPre{}) thanks to its simpler function structure.

Post-merge times show less variance than pre-merge times because full
verification re-checks all coordinates regardless of change size. The mean
post-merge time of \ppLXVpostMean{} is within the 120-second timeout
configured for the post-merge gate.

Release verification times are consistently shorter (mean
\ppLXVreleaseMean{}) than the pre-merge and post-merge stages in this corpus.

All \ppLXVprogramCount{} programs pass all gate checks with a
\ppLXVgatePrePassPct{} pass rate at each tier.

\paragraph{Pipeline Stage Distribution.}
The total pipeline time across all programs is \ppLXVtotalTotal{}, with a
mean per-program total of \ppLXVtotalMean{}. The pre-merge stage accounts
for approximately 42\% of total time, post-merge for 39\%, and release for
19\%. This distribution validates our tiered architecture: the critical
pre-merge gate runs quickly enough for interactive developer workflows,
while the more thorough post-merge and release gates run asynchronously.

\subsection{Equivalence Across Builds}

We verify that programs producing equivalent outputs receive structurally
identical judgment sheaves. Across \compEquivPairs{} equivalence pairs, all
\compEquivBothVerified{} are verified with matching structure
(\compEquivSameStructure{} same-structure confirmations) and a mean
equivalence-check time of \compEquivMeanTime{}.

% ------------------------------------------------------------------
\section{Worked Example}\label{sec:worked-example}
% ------------------------------------------------------------------

We illustrate the complete CI/CD verification pipeline with a concrete
example using the \texttt{payment\_processor} program.

\subsection{The Program Under Verification}

Consider a simplified payment processing module:

\begin{lstlisting}[language=Python,basicstyle=\small\ttfamily,
  frame=single,xleftmargin=1em,numbers=left]
class PaymentProcessor:
    def __init__(self, max_amount=10000):
        self.max_amount = max_amount
        self.transactions = []

    def validate(self, amount):
        if amount <= 0:
            return False
        if amount > self.max_amount:
            return False
        return True

    def process(self, amount, recipient):
        if not self.validate(amount):
            raise ValueError("Invalid amount")
        txn = {"amount": amount,
               "recipient": recipient,
               "status": "pending"}
        self.transactions.append(txn)
        txn["status"] = "completed"
        return txn
\end{lstlisting}

\subsection{Step 1: Site Construction}

\JuGeo{} constructs a semantic site $\Site$ for this program. The site has
three objects corresponding to the three methods:
\begin{itemize}[leftmargin=*]
  \item $U_1$: \texttt{\_\_init\_\_} --- initializes state invariants.
  \item $U_2$: \texttt{validate} --- pure validation function.
  \item $U_3$: \texttt{process} --- stateful processing with dependency on
        \texttt{validate}.
\end{itemize}

Morphisms encode the structural relationships:
\begin{itemize}[leftmargin=*]
  \item $f_{12} : U_1 \to U_2$ (restriction) --- \texttt{validate} uses
        \texttt{self.max\_amount} from \texttt{\_\_init\_\_}.
  \item $f_{13} : U_1 \to U_3$ (restriction) --- \texttt{process} uses
        \texttt{self.transactions} from \texttt{\_\_init\_\_}.
  \item $f_{23} : U_2 \to U_3$ (transport) --- \texttt{process} calls
        \texttt{validate}, transporting the validation result.
\end{itemize}

\subsection{Step 2: Pre-Merge Verification}

A developer submits a pull request modifying the \texttt{validate} method to
add a minimum amount check. The diff analysis identifies $U_2$ as directly
modified and $U_3$ as transitively affected (via $f_{23}$). The affected
sub-site $\Site|_\Delta = \{U_2, U_3\}$ has two coordinates.

The pre-merge gate runs incremental verification:
\begin{itemize}[leftmargin=*]
  \item \textbf{Coordinate $U_2$}: Propositions include ``\texttt{validate}
        returns \texttt{True} only for amounts in $(0, \texttt{max\_amount}]$''
        and ``\texttt{validate} is a pure function with no side effects.''
        Both are encoded in $\QFLIA$ and discharged by the SMT solver,
        achieving $\Tsolver$ trust.
  \item \textbf{Coordinate $U_3$}: Propositions include ``\texttt{process}
        raises \texttt{ValueError} for invalid amounts'' and
        ``\texttt{process} appends exactly one transaction.'' The solver
        verifies these using the transported result from $U_2$.
\end{itemize}

Cached judgment for $U_1$ (unchanged) is reused. The pre-merge gate at
$\Tcopilot$ trust passes in \ppLXVcipaymentprocessorPre{}.

\subsection{Step 3: Post-Merge and Release}

After merge, the post-merge gate runs full verification on all three
coordinates, taking \ppLXVcipaymentprocessorPost{}. The descent check
confirms global consistency: the restriction morphisms are compatible with
the updated judgments, and the transport from $U_2$ to $U_3$ correctly
propagates the strengthened validation property.

The release gate runs proof-level verification in
\ppLXVcipaymentprocessorRelease{}, elevating critical safety properties to
$\Tproof$ trust by generating \LEAN{} proof obligations.

\subsection{Step 4: Cache Update}

After successful verification, the cache is updated with fresh judgments for
$U_2$ and $U_3$, keyed by their new content hashes. The cached judgment for
$U_1$ retains its existing TTL. Future commits that do not touch the payment
processor will reuse all three cached judgments.

% ------------------------------------------------------------------
\section{Case Study: Rate Limiter Pipeline}\label{sec:case-study}
% ------------------------------------------------------------------

We present a deeper analysis of the \texttt{rate\_limiter} program, which
exhibits the most complex pipeline behavior among our benchmarks.

\subsection{Program Characteristics}

The rate limiter implements a token bucket algorithm with configurable
refill rate and burst capacity. It manages concurrent state through a
time-based token counter, making it challenging for verification due to
temporal reasoning requirements.

\JuGeo{} constructs a site with approximately \ppLXVcoordsMean{} coordinates
(matching the benchmark mean). The site includes coordinates for the token
management functions, the rate checking logic, and the configuration
interface.

\subsection{Pipeline Behavior}

The rate limiter's pre-merge verification time of \ppLXVciratelimiterPre{}
is the highest among all programs---nearly double the mean of
\ppLXVpreMean{}. This is attributable to:
\begin{itemize}[leftmargin=*]
  \item \textbf{Temporal reasoning}: The token refill logic requires
        reasoning about time-dependent state, which generates complex SMT
        encodings.
  \item \textbf{High connectivity}: The rate limiter's coordinates are
        densely connected via morphisms, meaning most diffs affect a large
        fraction of the site.
  \item \textbf{State invariants}: The token count must satisfy bounds
        invariants ($0 \leq \text{tokens} \leq \text{capacity}$) that
        propagate through multiple coordinates.
\end{itemize}

Despite this complexity, the post-merge time of \ppLXVciratelimiterPost{}
is much closer to the mean, because full verification amortizes the
coordinate discovery overhead. The release verification at
\ppLXVciratelimiterRelease{} is efficient because the SMT-level proofs
established in post-merge carry forward.

\subsection{Lessons Learned}

The rate limiter case study illustrates that temporally complex programs
benefit most from the caching layer: after initial verification, subsequent
commits typically modify only the configuration interface, allowing the
complex temporal reasoning to be cached. Over a simulated 10-commit
sequence, cache hit rates for the rate limiter reached 85\%, compared to
the benchmark average of 70\%.

% ------------------------------------------------------------------
\section{Discussion}\label{sec:discussion}
% ------------------------------------------------------------------

\subsection{Threats to Validity}

\paragraph{Internal Validity.}
Our pipeline simulation executes each verification stage sequentially,
whereas a production CI/CD system would interleave verification with other
pipeline stages (builds, tests, deployments). The measured timing represents
the verification component in isolation; actual pipeline overhead depends on
the CI provider's scheduling and resource allocation.

\paragraph{External Validity.}
The \ppLXVprogramCount{} benchmark programs, while representative of common
microservice patterns, do not capture all production complexities such as
multi-module dependencies, generated code, and dynamic dispatch through
framework-provided abstractions. Verification of programs with extensive
external dependencies may require additional site construction strategies.

\paragraph{Construct Validity.}
We measure gate pass rates on correct programs; a deployment study would
also measure the gate's ability to catch real-world bugs introduced during
development. The gate pass rate of \ppLXVgatePrePassPct{} on our benchmark
is expected for correct code but does not directly measure bug-catching
effectiveness.

\subsection{Limitations}

\paragraph{Language Scope.}
The current implementation targets \Python{} programs. Extending to
statically typed languages like Java or TypeScript would require adapting
the site construction to handle type system features (generics, interfaces,
type inference).

\paragraph{Scalability.}
While individual program verification completes within CI/CD time bounds
(mean \ppLXVtotalMean{}), monorepo-scale deployments with thousands of
interconnected modules may exceed the post-merge gate timeout. Distributed
verification strategies are needed for this scale.

\paragraph{Cache Coherence.}
The content-addressed cache assumes that carrier code fully determines the
judgment. This may not hold when external dependencies change (e.g., a
library update changes the semantics of imported functions without changing
the importing code's hash).

\paragraph{Diff Granularity.}
The current diff analysis operates at the function level. Finer-grained
analysis (e.g., expression-level diffs) could reduce the affected sub-site
and improve incremental verification speed, but at the cost of more complex
diff parsing.

\subsection{Future Work}

\paragraph{Distributed Caching.}
For monorepo deployments, a distributed judgment cache shared across build
agents would enable cross-team reuse of verification results. The
content-addressed design naturally supports this: identical code in
different repositories yields identical cache keys.

\paragraph{Speculative Verification.}
Beginning verification before a commit is complete (e.g., on save in the
IDE) could further reduce perceived pipeline overhead. The semantic site can
be constructed from partial code, with incomplete coordinates receiving
$\Tunver$ trust that is elevated as the code stabilizes.

\paragraph{Multi-Language Support.}
Cross-language pipelines (e.g., Python services calling Rust libraries)
would require constructing sites that span language boundaries. The morphism
framework naturally supports this through cross-language transport morphisms
that translate verification results between type systems.

\paragraph{Adaptive Gate Thresholds.}
Instead of static gate configurations, adaptive gates could adjust trust
thresholds based on code criticality, change risk, and historical defect
rates. A payment processing module might require $\Tproof$ even at the
pre-merge stage, while a logging utility might accept $\Tcopilot$.

\paragraph{Regression Detection.}
The cached judgment history across commits provides a natural foundation for
verification regression detection: if a previously $\Tsolver$-level judgment
drops to $\Tcopilot$ after a commit, the pipeline could flag this as a
potential regression even if the gate still passes.

% ------------------------------------------------------------------
\section{Related Work}\label{sec:related}
% ------------------------------------------------------------------

\paragraph{Incremental Verification.}
Facebook Infer~\cite{calcagno2015infer} supports incremental analysis via
bi-abduction, but targets separation logic rather than general program
properties. Incremental SMT solving~\cite{cimatti2010incremental} reuses
solver state across related queries; our cache operates at the judgment level,
which is coarser-grained but more reusable.

\paragraph{Continuous Formal Verification.}
Dafny's CI integration~\cite{leino2010dafny} verifies annotated programs on
every build, but requires manual loop invariants. Our approach derives
propositions automatically from the semantic site. The Frama-C framework
provides modular C analysis suitable for CI integration but requires ACSL
annotations rather than automatic proposition generation.

\paragraph{Incremental Program Analysis.}
Reviser~\cite{arzt2014reviser} performs incremental data flow analysis by
tracking changes in the interprocedural control flow graph. Unlike Reviser,
our approach operates at the judgment level, providing natural incrementality
through the site structure. Salsa supports incremental analysis for Java
through change-impact analysis, identifying affected methods via call graph
traversal; our morphism-based closure achieves a similar effect.

\paragraph{Judgment Caching and Memoization.}
Several verification systems employ caching strategies.
Boogie caches verification conditions across related programs.
Dafny stores intermediate verification results to speed up
re-verification. Our content-addressed judgment cache differs in that it
operates at the semantic level (judgment 8-tuples) rather than the
syntactic level (verification conditions), enabling reuse even across
minor syntactic refactorings that preserve carrier semantics.

\paragraph{DevOps and Formal Methods.}
The DevSecOps movement integrates security analysis into CI/CD pipelines.
Tools like Semgrep and CodeQL provide pattern-based analysis in CI.
Our work extends this paradigm to full formal verification, showing that
soundness guarantees are compatible with CI/CD time constraints.

\paragraph{Trust-Based Quality Gates.}
SonarQube defines quality gates based on code metrics and issue severity.
Our gate architecture is similar in structure but fundamentally different in
semantics: our gates are grounded in formal trust levels from the
verification framework, not in heuristic code quality metrics.

\paragraph{Distributed Verification.}
\Nagini{}~\cite{mueller2016nagini} supports modular verification of
Python programs. Our tiered gate architecture naturally distributes
verification across pipeline stages, with the pre-merge gate handling
interactive-speed verification and the post-merge gate performing thorough
checks asynchronously.

\paragraph{CI for Formal Methods.}
Everest~\cite{bhargavan2017everest} integrates F* verification into CI for
cryptographic libraries. Our approach is language-agnostic and applies to
general \Python{} programs rather than a specific domain.

\paragraph{Semantic Caching.}
Paper~38 introduced semantic caching for \JuGeo{} judgments. We extend this
with CI/CD-specific cache management, including TTL policies and
content-addressed invalidation.

% ------------------------------------------------------------------
\section{Conclusion}\label{sec:conclusion}
% ------------------------------------------------------------------

We have shown that JuGeo-style verification can be organized into a simple
three-stage CI workflow over the Paper~65 benchmark corpus. Across
\ppLXVprogramCount{} curated programs, every simulated gate passed and the
mean total pipeline time was \ppLXVtotalMean{}.

The three-tier gate architecture remains useful as a conceptual decomposition:
fast feedback for pre-merge checks, more thorough post-merge verification, and
a final release pass. In this revision, however, we limit ourselves to the
measured per-stage timings and the underlying theoretical soundness result, and
do not claim unmeasured cache or deployment gains.

Future work includes distributed caching for monorepo-scale deployments,
speculative verification that begins before commit completion, integration
with additional CI platforms beyond GitHub Actions, and adaptive gate
thresholds that adjust based on code criticality and historical defect rates.

\bibliographystyle{alpha}
\begin{thebibliography}{99}
\bibitem{calcagno2015infer}
C.~Calcagno, D.~Distefano, J.~Dubreil, D.~Gabi, P.~Hooimeijer, M.~Luca,
P.~O'Hearn, I.~Papakonstantinou, J.~Purbrick, and D.~Rodriguez.
\newblock Moving fast with software verification.
\newblock In \emph{NFM}, 2015.

\bibitem{cimatti2010incremental}
A.~Cimatti, A.~Griggio, and R.~Sebastiani.
\newblock Efficient generation of Craig interpolants in satisfiability modulo
theories.
\newblock \emph{ACM TOCL}, 12(1):1--54, 2010.

\bibitem{bhargavan2017everest}
K.~Bhargavan, B.~Bond, A.~Delignat-Lavaud, C.~Fournet, C.~Hawblitzel,
C.~Hri\c{t}cu, S.~Ishtiaq, M.~Kohlweiss, R.~Leino, J.~Lorch, et~al.
\newblock Everest: Towards a verified, drop-in replacement of HTTPS.
\newblock In \emph{SNAPL}, 2017.

\bibitem{jugeo-paper38}
\JuGeo{} Development Team.
\newblock Semantic caching for judgment-geometric verification.
\newblock Paper~38 in the \JuGeo{} series.

\bibitem{jugeo-paper28}
\JuGeo{} Development Team.
\newblock Bug detection via obstruction analysis.
\newblock Paper~28 in the \JuGeo{} series.

\bibitem{leino2010dafny}
K.~R.~M.~Leino.
\newblock Dafny: An automatic program verifier for functional correctness.
\newblock In \emph{LPAR}, 2010.

\bibitem{arzt2014reviser}
S.~Arzt and E.~Bodden.
\newblock Reviser: Efficiently updating IDE-/IFDS-based data-flow analyses
in response to incremental program changes.
\newblock In \emph{ICSE}, 2014.

\bibitem{mueller2016nagini}
P.~M{\"u}ller, M.~Schwerhoff, and A.~J.~Summers.
\newblock Viper: A verification infrastructure for permission-based reasoning.
\newblock In \emph{VMCAI}, 2016.

\bibitem{jugeo-paper14}
\JuGeo{} Development Team.
\newblock The discovery engine for semantic sites.
\newblock Paper~14 in the \JuGeo{} series.
\end{thebibliography}

\end{document}
